\beginsong{Les Champs Élysées}[by={Joe Dassin}]
\beginverse
\[C]Je m'baladais sur l'ave\[G]nue le c\[Am]{\oe}ur ouvert à l'incon\[C7]nu
\[F]J'avais envie de dire bon\[C]jour \[D7]à n'importe \[G]qui
\[C]N'importe qui et ce fut \[G]toi \[Am]je t'ai dit n'importe \[C7]quoi
\[F]Il suffisait de te par\[C]ler \[F]pour t'appri\[G]voi\[C]ser
\endverse
\beginchorus
\[C]Aux \[E7]Champs-Ély\[Am]sées, \[C7]
\[F]Aux \[C]Champs-Ély\[D7]sées, \[G]
\[C]Au soleil, sous la \[E7]pluie, \[Am]à midi ou à mi\[C7]nuit
\[F]Il y a tout ce que \[C]vous voulez aux \[F]Champs-Él\[G]ysé\[C]es
\endchorus
\beginverse
\[C]Tu m'as dit « J'ai rendez-\[G]vous \[Am]dans un sous-sol avec des \[C7]fous
\[F]Qui vivent la guitare à la \[C]main \[D7]du soir au ma\[G]tin »
\[C]Alors je t'ai accompa\[G]gnée \[Am]on a chanté, on a dan\[C7]sé
\[F]Et l'on n'a même pas pen\[C]sé \[F]à s'embras\[G]s\[C]er
\endverse
\beginverse
\[C]Hier soir deux in\[G]connus \[Am]et ce matin sur l'ave\[C7]nue
\[F]Deux amoureux tout étour\[C]dis \[D7]par la longue \[G]nuit
\[C]Et de l'Étoile à la Con\[G]corde \[Am]un orchestre à mille \[C7]cordes
\[F]Tous les oiseaux du point du \[C]jour \[F]chantent l'a\[G]m\[C]our
\endverse
\endsong

\beginsong{Santiano}[by={Hugues Aufray}]
\beginverse
\[Em]C'est un fameux trois-mâts fin comme un oi\[D]seau
\[Em]Hissez haut, Santi\[D]ano
\[Am]Dix-huit n{\oe}uds, quatre \[D]cent ton\[Bm]neaux
\[Em]Je suis fier d'y \[Bm]être mate\[Em]lot
\endverse
\beginchorus
\[Em]Tiens bon la vague et tiens bon le \[D]vent
\[Em]Hissez haut, Santi\[D]ano
\[Am]Si Dieu veut, toujours droit de\[D]van\[Bm]t
\[Em]Nous irons jus\[Bm]qu'à San Fran\[Em]cisco
\endchorus
\beginverse
\[Em]Je pars pour de longs mois en laissant Mar\[D]got
\[Em]Hissez haut, Santi\[D]ano
\[Am]D'y penser j'avais le c\[D]{\oe}ur \[Bm]gros
\[Em]En doublant les \[Bm]feux de Saint-Ma\[Em]lo
\endverse
\beginverse
\[Em]On prétend que là-bas l'argent coule à \[D]flots
\[Em]Hissez haut, Santi\[D]ano
\[Am]On trouve l'or au fond des ruis\[D]seau\[Bm]x
\[Em]J'en ramène\[Bm]rai plusieurs lin\[Em]gots
\endverse
\beginverse
\[Em]Un jour je reviendrai chargé de ca\[D]deaux
\[Em]Hissez haut, Santi\[D]ano
\[Am]Au pays j'irai voir Mar\[D]go\[Bm]t
\[Em]À son doigt je \[Bm]passerai l'an\[Em]neau
\endverse
\beginchorus
\[Em]Tiens bon le cap et tiens bon le \[D]flot
\[Em]Hissez haut, Santi\[D]ano
\[Am]Sur la mer qui fait le gros \[D]do\[Bm]s
\[Em]Nous irons jus\[Bm]qu'à San Fran\[Em]cisco
\endchorus
\endsong

\beginsong{Cendrillon}[by={Téléphone}]
\beginverse
\[A]Cendrillon, pour \[E]ses vingt ans
Est \[F#m]la plus jolie \[D]des enfants
Son \[A]bel amant le \[E]Prince charmant
La \[F#m]prend sur son \[D]cheval blanc
Elle \[E]oublie le \[A]temps dans \[E]ce palais d'ar\[F#m]gent
Pour \[Bm]ne pas voir qu'un nouveau jour se lève
Elle \[D]ferme les yeux et dans ses rêves
\endverse
\beginchorus
Elle \[A]part \[E] \[F#m]
\[C#m]Jolie petite \[D]his\[A]toire \[E] \[F#m] \[C#m] \[D]
Elle \[A]part \[E] \[F#m]
\[C#m]Jolie petite \[D]his\[A]toire \[E] \[F#m] \[C#m] \[D]
\endchorus
\beginverse
\[A]Cendrillon, pour \[E]ses trente ans
Est \[F#m]la plus triste \[D]des mamans
Le \[A]Prince charmant a \[E]foutu le camp
A\[F#m]vec la belle au \[D]bois dormant
\[A]Elle a vu cent \[E]chevaux blancs
\[F#m]Loin d'elle emmener \[D]ses enfants
\[E]Elle commence à \[A]boire à traî\[E]ner dans les \[F#m]bars
Em\[Bm]mitouflée dans son cafard
Main\[D]tenant elle fait le trottoir
\endverse
\beginchorus
Elle \[A]part \[E] \[F#m]
\[C#m]Jolie petite \[D]his\[A]toire \[E] \[F#m] \[C#m] \[D]
Elle \[A]part \[E] \[F#m]
\[C#m]Jolie petite \[D]his\[A]toire \[E] \[F#m] \[C#m] \[D]
\endchorus
\beginverse
\[A]Dix ans de cette vie \[E]ont suffi
À \[F#m]la changer en \[D]junkie
Et \[A]dans un sommeil \[E]infini
\[F#m]Cendrillon voit fi\[D]nir sa vie
\[E]Les lumières \[A]dansent \[E]dans l'ambu\[F#m]lance
Et \[Bm]elle tue sa dernière chance
\[D]Tout ça n'a plus d'importance
\endverse
\beginchorus
Elle \[A]part \[E] \[F#m]
\[C#m]Fin de l'\[D]his\[A]toire \[E] \[F#m]
\endchorus
\beginverse
\[A]Notre père qui \[E]êtes si vieux
As-\[F#m]tu vraiment fait \[D]de ton mieux ?
Car \[A]sur la terre, et \[E]dans les cieux
Tes \[F#m]anges n'aiment \[Dm]pas devenir \[A]vieux
\endverse
\endsong

\beginsong{Qui peut faire de la voile sans vent}
\beginverse
\[Am]Qui peut faire de la voile sans vent ?
\[Dm]Qui peut ramer sans \[Am]rames ?
\[Dm]Et qui peut quitter son \[Am]ami
\[E7]Sans verser une \[Am]larme ?
\endverse
\beginverse
\[Am]Je peux faire de la voile sans vent,
\[Dm]Je peux ramer sans \[Am]rames,
\[Dm]Mais ne peux quitter mon \[Am]ami
\[E7]Sans verser une \[Am]larme
\endverse
\beginverse
\[Am]Qui peut voir se coucher le Soleil
\[Dm]Sans que la nuit ne \[Am]tombe ?
\[Dm]Et qui peut retrouver le som\[Am]meil
\[E7]Sans que son c{\oe}ur s'ef\[Am]fondre ?
\endverse
\beginverse
\[Am]Je peux voir se coucher le soleil
\[Dm]Sans que la nuit ne \[Am]tombe
\[Dm]Mais je ne peux trouver le som\[Am]meil
\[E7]Sans que mon c{\oe}ur s'ef\[Am]fondre
\endverse
\beginverse
\[Am]Qui peut trouver l'oiseau sans ramage ?
\[Dm]Qui peut vivre sans ri\[Am]vage ?
\[Dm]Mais qui peut s'accrocher au ro\[Am]cher
\[E7]Sans faire sombrer une ami\[Am]tié ?
\endverse
\beginverse
\[Am]Qui peut croire un instant à l'amour
\[Dm]Quand tant d'hommes se \[Am]battent ?
\[Dm]Et qui ne peut oublier un \[Am]jour
\[E7]Le monde et son mas\[Am]sacre ?
\endverse
\endsong

\beginsong{Ne sens-tu pas claquer tes doigts}
\beginverse
\[C]Ne sens-tu pas claquer tes doigts, claquer tes doigts
\[Am]Et la musique monter en toi, monter en toi
\[F]N'attends pas que le feu soit mort, le feu soit mort
\[G]Chante tant que tu peux encore, tu peux encore
\[C]Encore, encore, encore, encore, encore.
\endverse
\beginverse
\[C]Tu peux donner tout c'que tu as, tout c'que tu as,
\[Am]Mais oui la vie c'est fait pour ça, c'est fait pour ça.
\[F]Tu peux sourire autour de toi, autour de toi,
\[G]Tendre la main à qui voudra, à qui voudra.
\[C]La main, la main, la main, la main, la main.
\endverse
\beginverse
\[C]Ne sens-tu pas chanter la vie, chanter la vie,
\[Am]Et toute la joie réunis, joie réunis,
\[F]Qui nous rassemble tous ici, oui tous ici,
\[G]Jusqu'à la fin de notre vie.
\[C]La vie, la vie, la vie, la vie, la vie.
\endverse
\beginverse
\[C]Et pour tous ceux qui ont faim, ceux qui ont faim,
\[Am]Donne-leur un peu de ton pain, peu de ton pain,
\[F]Pour qu'ils partagent tous ensemble, oui tous ensemble,
\[G]Cet amour qui nous rassemble, oui nous rassemble.
\[C]L'amour, l'amour, l'amour, l'amour, l'amour.
\endverse
\beginverse
\[C]Ne sens-tu pas battre ton c{\oe}ur, battre ton c{\oe}ur,
\[Am]Qui s'éparpille en mille fleurs, en mille fleurs.
\[F]Et prends la main se ton ami, de ton ami,
\[G]Et garde la toute ta vie, toute ta vie.
\[C]La vie, la vie, la vie, la vie, la vie.
\endverse
\endsong

\beginsong{Debout les gars}[by={Hugues Aufray}]
\beginverse
\[Am]Cette montagne que tu vois
\[G]On en viendra à bout, mon \[C]gars
\[Am]Un bulldozer et deux cent bras
\[G]Et passera la \[Am]route
\endverse
\beginchorus
\[Am]Debout les gars ! Réveillez-vous !
\[G]Y va falloir en mettre un \[C]coup
\[Am]Debout les gars ! Réveillez-vous !
\[G]On va au bout du \[Am]monde
\endchorus
\beginverse
\[Am]Il ne faut pas se dégonfler
\[G]Devant des tonnes de ro\[C]cher
\[Am]On va faire un quatorze juillet
\[G]À coups de dyna\[Am]mite
\endverse
\beginverse
\[Am]Encore un mètre et deux et trois
\[G]En mille neuf cent quatre-vingt-\[C]trois
\[Am]Tes enfants seront fiers de toi
\[G]La route sera \[Am]belle
\endverse
\beginverse
\[A#m]Il nous arrive, parfois, le soir
\[G#]Comme un petit coup de ca\[C#]fard
\[A#m]Mais ce n'est qu'un peu de brouillard
\[G#]Que le soleil dé\[A#m]chire
\endverse
\beginverse
\[A#m]Les gens nous prenaient pour des fous
\[G#]Mais nous, on passera par\[C#]tout
\[A#m]Et nous serons au rendez-vous
\[G#]De ceux qui nous at\[A#m]tendent
\endverse
\beginverse
\[A#m]Et quand tout sera terminé
\[G#]Y faudra bien se sépa\[C#]rer
\[A#m]Mais on oubliera jamais
\[G#]Ce qu'on a fait en\[A#m]semble
\endverse
\endsong

\beginsong{Le petit âne gris}[by={Hugues Aufray}]
\beginverse
\[Em]Écoutez cette histoire
\[D]Que l'on m'a racon\[Em]tée
Du fond de ma mémoire,
\[D]Je vais vous la chan\[Em]ter
\[G]Elle se passe en Pro\[D]vence,
\[C]Au milieu des mou\[D]tons,
\[Em]Dans le sud de la France,
\[G]Au pa\[D]ys des san\[Em]tons
\endverse
\beginverse
\[Em]Quand il vint au domaine,
\[D]Y'avait un beau trou\[Em]peau
Les étables étaient pleines
\[D]De brebis et d'a\[Em]gneaux
\[G]Marchant toujours en \[D]tête
\[C]Aux premières lu\[D]eurs,
\[Em]Pour tirer sa charrette,
\[G]Il met\[D]tait tout son \[Em]c{\oe}ur
\endverse
\beginverse
\[Em]Au temps des transhumances,
\[D]Il s'en allait heu\[Em]reux,
Remontant la Durance,
\[D]Honnête et coura\[Em]geux
\[G]Mais un jour, de Mar\[D]seille,
\[C]Des messieurs sont ve\[D]nus
\[Em]La ferme était bien vieille,
\[G]Alors \[D]on l'a ven\[Em]due
\endverse
\beginverse
\[Em]Il resta au village
\[D]Tout le monde l'aimait \[Em]bien,
Vaillant, malgré son âge
\[D]Et malgré son cha\[Em]grin
\[G]Image d'évan\[D]gile,
\[C]Vivant d'humili\[D]té,
\[Em]Il se rendait utile
\[G]Auprès \[D]du canton\[Em]nier
\endverse
\beginverse
\[Em]Cette vie honorable,
\[D]Un soir, s'est termi\[Em]née
Dans le fond d'une étable,
\[D]Tout seul il s'est cou\[Em]ché
\[G]Pauvre bête de \[D]somme,
\[C]Il a fermé les \[D]yeux
\[Em]Abandonné des hommes,
\[G]Il est \[D]mort sans a\[Em]dieux
\endverse
\beginverse
Mm mm mmm mm...
\[G]Cette chanson sans \[D]gloire
\[C]Vous racontait la \[D]vie,
\[Em]Vous racontait l'histoire
\[G]D'un pe\[D]tit âne \[Em]gris...
\endverse
\endsong

\beginsong{Je cherche fortune}[by={Aristide Bruant}]
\beginverse
\[F]Chez l'boulanger (bis)
Fais-moi crédit (bis)
\[C]J'n'ai pas d'argent (bis)
\[F]J'paierai sam'di (bis)
Si tu n'veux pas (bis)
M'donner du pain (bis)
J'te fourre la tête (bis)
\[F]Dans ton pétrin (bis)
\endverse
\beginverse
Non, c'est pas moi, c'est ma s{\oe}ur
Qu'a cassé la machine à vapeur
Ta gueule (ter)
\endverse
\beginchorus
\[F]Je cherche fortune !
\[C]Autour du Chat Noir
\[Bb]Au clair de la \[F]lune
\[C7]À Montmartre, le \[F]soir
\endchorus
\beginverse
\[F]Chez l'marchand d'frites...
M'donner des frites
J'te cass' la gueule
Dans tes marmites
\endverse
\beginverse
\[F]Chez l'cabar'tier...
M'donner à boire
J'te cass' la gueule
Sur ton comptoir
\endverse
\beginverse
\[F]Marchand d'tabac...
M'donner des sèches
J'fais dans ta gueule
Un'large brèche
\endverse
\beginverse
\[F]Chez l'aubergiste...
M'donner un'chambre
J'te cass' la gueule
Et les cinq membres
\endverse
\beginverse
\[F]Chez l'pharmacien...
M'donner d'potion
J'te cass' la gueule
Dans tes flacons
\endverse
\beginverse
\[F]Chez M'sieur l'curé...
Nous mari-er
J'te cass' la gueule
Dans l'bénitier
\endverse
\endsong

\end{songs}\songcolumns{2}\begin{songs}{titleidx}

\beginsong{Armstrong}[by={Claude Nougaro}]
\beginverse
\[Em]Armstrong, je ne \[B7]suis pas \[Em]noir, \[Am]je \[Em]suis \[B7]blanc de \[Em]peau \[A7]
\[Em]Quand on veut chan\[B7]ter l'es\[Em]poir \[Am]quel \[Em]manque de \[B7]pot ! \[Em]
\[Em]Oui, j'ai beau voir le \[Am]ciel, l'oi\[Em]seau, rien rien \[Em/D]rien ne luit là-\[Cmaj7]haut \[Am]
\[Em]Les anges, \[A7]zéro, \[Em]je suis \[B7]blanc de \[Em]peau
\endverse
\beginverse
\[Em]Armstrong, tu te \[B7]fends la \[Em]poire, \[Am]on \[Em]voit toutes tes \[B7]dents \[Em] \[A7]
\[Em]Moi, je broie plu\[B7]tôt du \[Em]noir, \[Am]du \[Em]noir en de\[B7]dans \[Em]
\[Em]Chante pour moi, \[Am]Louis, oh \[Em]oui, chante chante \[Em/D]chante, ça tient \[Cmaj7]chaud \[Am]
\[Em]J'ai froid, oh \[A7]moi, \[Em]qui suis \[B7]blanc de \[Em]peau
\endverse
\beginverse
\[Em]Armstrong, la vie, quelle his\[B7]toire, \[Em]c'est pas très mar\[Am]rant \[Em] \[B7] \[Em] \[A7]
\[Em]Qu'on l'écrive blanc sur \[B7]noir \[Em]ou bien \[Am]noir sur \[Em]blanc \[B7] \[Em]
\[Em]On voit surtout du \[Am]rouge, du \[Em]rouge, sang sang sans \[Em/D]trêve ni re\[Cmaj7]pos \[Am]
\[Em]Qu'on soit, ma \[A7]foi, \[Em]noir ou \[B7]blanc de \[Em]peau
\endverse
\beginverse
\[Em]Armstrong, un jour, tôt ou \[B7]tard, \[Em]on n'est que des \[Am]os \[Em] \[B7] \[Em] \[A7]
\[Em]Est-ce que les tiens se\[B7]ront \[Em]noirs ? Ce s'\[Am]rait rigo\[Em]lo \[B7] \[Em]
\[Em]Allez, Louis, allé\[Am]luia, \[Em]au-de\[B7]là de nos ori\[G7]peaux \[A7]
\[Em]Noir et blanc sont ressem\[B7]blants comme deux gouttes d'\[Em]eau
\[A]Oh \[Em]Yeah
\endverse
\endsong

\end{songs}\songcolumns{3}\begin{songs}{titleidx}

\beginsong{Chevaliers de la Table Ronde}
\beginverse
\[C]Chevaliers de la Table Ronde
\[G7]Goûtons voir si le vin est \[C]bon
\[C]Chevaliers de la Table Ronde
\[G7]Goûtons voir si le vin est \[C]bon \[C7]
\endverse
\beginchorus
\[F]Goûtons voir, oui, oui, oui
\[C]Goûtons voir, non, non, non
\[G7]Goûtons voir si le vin est \[C]bon \[C7]
\[F]Goûtons voir, oui, oui, oui
\[C]Goûtons voir, non, non, non
\[G7]Goûtons voir si le vin est \[C]bon
\endchorus
\beginverse
\[C]J'en boirai cinq à six bouteilles
\[G7]Une femme sur mes ge\[C]noux
\[C]J'en boirai cinq à six bouteilles
\[G7]Une femme sur mes ge\[C]noux \[C7]
\endverse
\beginverse
\[C]Si je meurs je veux qu'on m'enterre
\[G7]Dans une cave où y'a du bon \[C]vin
\[C]Si je meurs je veux qu'on m'enterre
\[G7]Dans une cave où y'a du bon \[C]vin \[C7]
\endverse
\beginverse
\[C]Les deux pieds contre la muraille
\[G7]Et la tête sous le robi\[C]net
\[C]Les deux pieds contre la muraille
\[G7]Et la tête sous le robi\[C]net \[C7]
\endverse
\beginverse
\[C]Sur ma tombe je veux qu'on inscrive
\[G7]Ici gît le roi des bu\[C]veurs
\[C]Sur ma tombe je veux qu'on inscrive
\[G7]Ici gît le roi des bu\[C]veurs
\endverse
\endsong

\beginsong{Comme un homme}[by={Mulan}]
\beginverse
\[Em]Attaquons l'exer\[D/F#]cice pour \[G]défaire les \[D/F#]Huns, \[G]
\[Am]M'ont-ils donné leurs fils ? Je n'en \[D]vois pas un.
\[C/E]Vous êtes plus fragiles que des \[D/F#]fillettes,
Mais \[G]jusqu'au bout, et \[C]coup par coup,
Je \[C]saurai \[Dsus4]faire \[D]de vrais hommes de \[Em]vous. \[D] \[Em]
\endverse
\beginverse
\[Em]Comme la flèche qui \[D/F#]vibre, et \[G]frappe en plein \[D/F#]c{\oe}ur, \[G]
\[Am]En trouvant l'équilibre, vous serez vain\[D]queur.
\[C/E]Vous n'êtes qu'une bande de \[D/F#]femmelettes,
Mais \[G]envers et contre \[C]tout,
Je \[C]saurai \[Dsus4]faire \[D]de vrais hommes de \[Em]vous ! \[D] \[Em]
\endverse
\beginverse*
\[C]J'aurai dû me mettre au ré\[D]gime...
\[B]Saluez tous mes amis pour \[Em]moi...
\[D/F#]Je n'aurai pas dû sécher les \[G]cours de \[C]gym...
\[C]Ce gars-là nous flanque les \[D]foies...
\[B]Mais s'il voyait la fille en \[Em]moi ?
\[D/F#]Je suis tout en nage mais \[G]nager je ne sais \[C]pas...
\endverse*
\beginchorus
\[D](Comme un \[C/E]homme !)
\[C]Sois plus violent que le \[D]cours du tor\[B]rent. \[Em]
\[D](Comme un \[C/E]homme !)
\[C]Sois plus puissant que les \[D]oura\[B]gans. \[Em]
\[C](Comme un homme !)
\[C]Sois plus ardent que le \[D]feu des vol\[B]cans. \[Em]
\[C]Secret comme les nuits de \[D]lune de l'O\[Esus4]ri\[E]ent.
\endchorus
\beginverse
\[F#m]Les jours passent et les \[E/G#]Huns
Ne \[A]sont plus, \[E/G#]très \[A]loin. \[Bm] \[E]
\[F#m]Suivez bien mon \[E/G#]chemin,
\[A]Vous vi\[E/G#]vrez de\[A]main. \[Bm] \[E]
\[D]Vous ne serez jamais vaillants et \[E]forts,
Comme des \[A]hommes, rentrez chez \[D]vous,
Je \[D]ne peux faire de vrais \[E]hommes de \[F#m]vous.
\endverse
\endsong

\beginsong{Il en faut peu pour être heureux}[by={Le Livre de la Jungle}]
\beginverse
\[G]Il en faut \[G7]peu pour être heu\[C]reux,
Vraiment très peu pour être heu\[C7]reux.
\[G]Il faut se satis\[E7]faire du néces\[A7]saire, \[D7]
\[G]Un peu d'eau fraîche et de ver\[G7]dure
Que nous \[C]prodigue la na\[C7]ture,
\[G]Quelques \[E7]rayons de \[A7]miel et de so\[D7]leil. \[G]
\endverse
\beginverse
\[D7]Je dors d'ordinaire sous les fron\[G]daisons,
Et \[D7]toute la jungle est ma mai\[G]son. \[G7]
\[C]Toutes les abeilles de la fo\[Cm]rêt
Butinent \[G]pour moi dans les bos\[A7]quets.
\[A7]Et quand je retourne un gros caillou,
\[D7]Je sais trouver des fourmis des\[G]sous,
Es\[E7]saie, c'est bon, c'est \[A7]doux !
Il en faut \[D7]vraiment peu, très peu pour être heu\[G]reux.
\endverse
\beginchorus
\[G]Il en faut \[G7]peu pour être heu\[C]reux,
Vraiment très peu pour être heu\[C7]reux.
\[G]Chassez de votre es\[E7]prit tous vos sou\[A7]cis, \[D7]
\[G]Prenez la vie du bon cô\[G7]té,
Riez, \[C]sautez, dansez, chan\[C7]tez,
Et \[G]vous \[E7]serez un \[A7]ours très \[D7]bien lé\[G]ché.
\endchorus
\beginverse
\[D7]Cueillir une banane, oui ! Ça se fait sans as\[G]tuce,
Mais \[D7]c'est tout un drame si c'est un cac\[G]tus. \[G7]
\[C]Si vous chipez des fruits sans é\[Cm]pine,
Ce n'est \[G]pas la peine de faire atten\[A7]tion.
\[A7]Mais si le fruit de vos rapines est tout plein d'épines,
\[D7]C'est beaucoup moins \[G]bon.
A\[E7]lors petit, as-tu com\[A7]pris ?
Il en faut \[D7]vraiment peu, très peu pour être heu\[G]reux.
\endverse
\beginverse*
\[E7]Et tu verras qu'tout est résolu
\[A7]Lorsque l'on se passe
\[D7]Des choses superflues,
\[G]Alors tu ne t'en \[E7]fais plus.
Il en faut \[A7]vraiment peu, très \[D7]peu pour être heu\[G]reux.
\endverse*
\endsong

\beginsong{J't'emmène au vent}[by={Louise Attaque}]
\textnote{Capo 4}
\beginchorus
\[Em]Allez viens j't'emmène au \[G]vent,
\[Em]Je t'emmène au-dessus des \[G]gens.
\[Am]Et je voudrais que tu te rap\[Em]pelles,
Notre amour est éter\[D]nel et pas artifi\[F]ciel.
\endchorus
\beginverse
\[Em]Je voudrais que tu te ramènes de\[G]vant,
\[Em]Que tu sois là de temps en \[G]temps.
\[Am]Et je voudrais que tu te rap\[Em]pelles,
Notre amour est éter\[D]nel et pas artifi\[F]ciel.
\endverse
\beginverse
\[Em]Je voudrais que tu m'appelles plus sou\[G]vent,
\[Em]Que tu prennes parfois l'de\[G]vant.
\[Am]Et je voudrais que tu te rap\[Em]pelles,
Notre amour est éter\[D]nel et pas artifi\[F]ciel.
\endverse
\beginverse
\[Em]Je voudrais que tu sois celle que j'en\[G]tends,
\[Em]Allez viens j't'emmène au-dessus des \[G]gens.
\[Am]Et je voudrais que tu te rap\[Em]pelles,
Notre amour est éter\[D]nel, artifi\[F]ciel.
\endverse
\endsong

\beginsong{L'hymne de nos campagnes}[by={Tryo}]
\beginverse
\[Dm]Si tu es né dans une cité \[Bb]HLM, \[A]
\[Dm]Je te dédicace ce po\[Bb]ème. \[A]
\[Dm]Eh, les mans, faut faire la \[Bb]part des \[A]choses,
\[Dm]Il est grand temps de faire une \[Bb]pause. \[A]
\[Dm]De troquer cette vie mo\[Bb]rose \[A]
\[Dm]Contre le parfum d'une \[Bb]rose. \[A]
\endverse
\beginchorus
\[Dm]C'est l'hymne de nos cam\[Bb]pagnes, \[A]
\[Dm]De nos rivières, de nos mon\[Bb]tagnes, \[A]
\[Dm]De la vie man, du monde ani\[Bb]mal, \[A]
\[Dm]Crie-le bien fort, use tes cordes vo\[Bb]cales ! \[A]
\endchorus
\beginverse
\[Dm]Pas de boulot, pas de di\[Bb]plôme, \[A]
\[Dm]Partout la même odeur de \[Bb]zone. \[A]
\[Dm]Va voir ailleurs, rien ne te re\[Bb]tient, \[A]
\[Dm]Va vite faire quelque chose de tes \[Bb]mains. \[A]
\[Dm]Ne te retourne pas, ici tu n'as \[Bb]rien, \[A]
\[Dm]Et sois le premier à chanter ce re\[Bb]frain. \[A]
\endverse
\beginverse
\[Dm]Assieds-toi près d'une ri\[Bb]vière, \[A]
\[Dm]Écoute le coulis de l'eau sur la \[Bb]terre. \[A]
\[Dm]Dis-toi qu'au bout, hé ! il y a la \[Bb]mer, \[A]
\[Dm]Et que ça, ça n'a rien d'éphé\[Bb]mère. \[A]
\endverse
\beginverse
\[Dm]Assieds-toi près d'un vieux \[Bb]chêne \[A]
\[Dm]Et compare-le à la race hu\[Bb]maine. \[A]
\[Dm]Lève la tête, regarde ses \[Bb]feuilles, \[A]
\[Dm]Tu verras peut-être un écu\[Bb]reuil. \[A]
\endverse
\beginverse
\[Dm]Peut-être que je parle pour ne rien \[Bb]dire, \[A]
\[Dm]Que quand tu m'écoutes tu as envie de \[Bb]rire. \[A]
\[Dm]Et j'aimerais pour tous les ani\[Bb]maux \[A]
\[Dm]Que tu captes le message de mes \[Bb]mots. \[A]
\[Dm]Car un lopin de terre, une tige de ro\[Bb]seau \[A]
\[Dm]Servira à la croissance de tes mar\[Bb]mots ! \[A]
\endverse
\endsong

\beginsong{Le lion est mort ce soir}[by={Pow Wow}]
\beginverse*
\[G]Wé-ah, wé-\[C]ah, \[G]wé-ah, wé-\[D]ah.
\endverse*
\beginverse
\[G]Dans la jungle, terrible \[C]jungle,
Le \[G]lion est mort ce \[D]soir.
\[G]Et les hommes tranquilles s'en\[C]dorment,
Le \[G]lion est mort ce \[D]soir.
\endverse
\beginchorus
\[G]Owimboe, o wimboe, owimboe, o \[C]wimboe,
\[G]Owimboe, o wimboe, owim\[D]boe.
\endchorus
\beginverse
\[G]Tout est sage dans le vil\[C]lage,
Le \[G]lion est mort ce \[D]soir.
\[G]Plus de rage, plus de car\[C]nage,
Le \[G]lion est mort ce \[D]soir.
\endverse
\beginverse
\[G]L'indomptable, le redou\[C]table,
Le \[G]lion est mort ce \[D]soir.
\[G]Viens ma belle, viens ma ga\[C]zelle,
Le \[G]lion est mort ce \[D]soir.
\endverse
\endsong

\beginsong{Le matou revient}[by={Steve Waring}]
\beginverse
\[Am]Tompson, le vieux fer\[G]mier, a beau\[F]coup d'en\[E]nuis.
\[Am]Il n'arrive pas à se \[G]débarrasser de son \[F]vieux gros chat \[E]gris.
\[Am]Pour mettre à la porte son \[G]chat, il a ten\[F]té n'importe \[E]quoi.
\[Am]Il l'a même posté au Ca\[G]nada et lui a \[F]dit « Tu resteras \[E]là ! »
\endverse
\beginchorus
\[Am]Mais le matou re\[G]vient le jour sui\[F]vant,
Le matou re\[E]vient, il est toujours vivant.
\endchorus
\beginverse
\[Am]Tompson paie un petit \[G]gars pour assas\[F]siner le \[E]chat.
\[Am]L'enfant part à la \[G]pêche, l'animal \[F]dans les \[E]bras.
\[Am]Au milieu de la ri\[G]vière, le ca\[F]not a cou\[E]lé.
\[Am]Le fermier apprend que \[G]l'enfant s'est \[F]noyé\[E]ééé.
\endverse
\beginverse
\[Am]Le voisin de Tompson com\[G]mence à s'éner\[F]ver,
Il prend sa cara\[E]bine et la bourre de T.N.T.
\[Am]Le fusil éclate, la \[G]ville est affo\[F]lée,
Une pluie de petits \[E]morceaux d'homme vient de tomber.
\endverse
\beginverse
\[Am]Le fermier décou\[G]ragé envoie son \[F]chat chez le bou\[E]cher
\[Am]Pour qu'il en fasse du \[G]hachis Parmen\[F]tier.
Le chat hurle et dis\[E]paraît dans la machine.
\[Am]« De la viande poi\[G]lue » est affi\[F]ché sur la vi\[E]trine.
\endverse
\beginverse
\[Am]Un fou s'engage à \[G]partir en bal\[F]lon
Pour aller dans la \[E]lune déposer le chaton.
\[Am]Au cours du voyage, le \[G]ballon a cre\[F]vé.
À l'autre bout du \[E]monde, un cadavre est retrouvé.
\endverse
\beginverse
\[Am]Cette fois-ci, on en\[G]voie le chat au \[F]Cap Kenne\[E]dy.
\[Am]C'est dans une fusée à \[G]trois étages \[F]qu'il est par\[E]ti.
\[Am]Le fermier saute de \[G]joie, car il n'a \[F]plus de sou\[E]cis.
\[Am]Le lendemain matin, \[G]on l'appelle \[F]de Mia\[E]mi...
\endverse
\endsong

\beginsong{Nagawicka}[by={Jacky Galou}]
\beginverse
\[Dm]Un petit indien, un petit indien,
Na\[A]gawic\[Dm]ka, Na\[A]gawic\[Dm]ka.
\[F]Chantait gaie\[Gm]ment sur le che\[C]min, \[F]
Na\[Gm]gawic\[Dm]ka, Na\[A]gawic\[Dm]ka.
\endverse
\beginverse
\[Dm]Quand je serai grand, quand je serai grand,
Na\[A]gawic\[Dm]ka, Na\[A]gawic\[Dm]ka.
\[F]J'aurai un \[Gm]arc et un car\[C]quois, \[F]
Na\[Gm]gawic\[Dm]ka, Na\[A]gawic\[Dm]ka.
\endverse
\beginverse
\[Dm]Avec mes flèches, avec mes flèches,
Na\[A]gawic\[Dm]ka, Na\[A]gawic\[Dm]ka.
\[F]Je chasse\[Gm]rai le grand bi\[C]son, \[F]
Na\[Gm]gawic\[Dm]ka, Na\[A]gawic\[Dm]ka.
\endverse
\beginverse
\[Dm]Sur mon cheval, sur mon cheval,
Na\[A]gawic\[Dm]ka, Na\[A]gawic\[Dm]ka.
\[F]J'irai plus \[Gm]vite que le \[C]vent, \[F]
Na\[Gm]gawic\[Dm]ka, Na\[A]gawic\[Dm]ka.
\endverse
\beginverse
\[Dm]Autour du feu, autour du feu,
Na\[A]gawic\[Dm]ka, Na\[A]gawic\[Dm]ka.
\[F]Je danse\[Gm]rai toute la \[C]nuit, \[F]
Na\[Gm]gawic\[Dm]ka, Na\[A]gawic\[Dm]ka.
\endverse
\endsong

\beginsong{Partenaire particulier}[by={Partenaire Particulier}]
\beginverse
\[F#m]Je suis un \[D]être, à la re\[E]cherche, \[C#m]
\[F#m]Non pas de la vé\[D]rité, \[E] \[C#m]
\[F#m]Mais simple\[D]ment d'une aven\[E]ture \[C#m]
\[F#m]Qui sorte un peu \[D]de la banali\[E]té. \[C#m]
\endverse
\beginverse
\[F#m]J'en ai as\[D]sez de ce car\[E]can, \[C#m]
\[F#m]Qui m'enferme \[D]dans toutes ses \[E]règles. \[C#m]
\[F#m]Il me dit \[D]de rester dans la \[E]norme, \[C#m]
\[F#m]Mais l'on finit \[D]par s'y ennu\[E]yer. \[C#m]
\endverse
\beginverse*
\[D]Alors je cherche et je \[E]trouverai
Cette \[A]fille qui me manque \[F#m]tant.
\[D]Alors je cherche et je \[E]trouverai
Cette \[A]fille qui me tente \[F#m]tant,
Qui me \[D]tente \[E]tant !
\endverse*
\beginchorus
\[C#m/E]Partenaire parti\[F#m]culier cherche \[D]partenaire particu\[E]lière,
\[C#m/E]Débloquée, pas \[F#m]trop timide \[D]et une bonne dose de savoir-\[E]faire.
\endchorus
\beginverse
\[F#m]Vous compren\[D]drez que de tels pé\[E]chés \[C#m]
\[F#m]Sont parfois diffi\[D]ciles à a\[E]vouer. \[C#m]
\[F#m]Ils sont au\[D]tour de moi, si coin\[E]cés, \[C#m]
\[F#m]Ce n'est pas par\[D]mi eux que je \[E]trouverai. \[C#m]
\endverse
\endsong

\beginsong{S.O.S. d'un terrien en détresse}[by={Daniel Balavoine}]
\beginchorus
\[Am]Pourquoi je vis, pour\[Dm]quoi je \[Am]meurs,
\[F]Pourquoi je ris, pourquoi je \[E]pleure.
\[Am]Voici le \[Dm]S.O.S. d'un terrien en dé\[E]tresse.
\endchorus
\beginverse
\[Am]J'ai jamais eu les pieds sur \[F]terre,
J'aim'rais \[C]mieux être un oi\[Dm]seau,
J'suis mal dans ma \[Am]peau.
\[Am]J'voudrais voir le monde à l'en\[F]vers,
Si jamais \[C]c'était plus \[Dm]beau,
Plus beau vu d'en \[Am]haut.
\[F]Oh ohh... \[E]d'en haut.
\endverse
\beginverse
\[Am]J'ai toujours \[F]confondu la \[E]vie
\[Am]Avec \[F]les bandes dessi\[E]nées.
\[Am]J'ai comme des \[F]envies de méta\[C]morphose,
Je sens quelque \[E]chose
Qui m'at\[Am]tire, qui m'\[G]attire,
Qui m'at\[F]tire vers le \[E]haut.
\endverse
\beginverse
\[Am]Au grand loto de l'uni\[F]vers,
J'ai pas \[C]tiré le bon nu\[Dm]méro,
J'suis mal dans ma \[Am]peau.
\[F]Oh ohh... j'ai pas envie d'être un ro\[C]bot, \[G]
Métro boulot \[Am]dodo.
\endverse
\endsong

\beginsong{Toi plus moi}[by={Grégoire}]
\beginchorus
\[Am]Toi, plus moi, plus tous ceux qui le \[F]veulent,
Plus lui, plus elle et tous ceux qui sont \[C]seuls.
Allez, venez et entrez dans la \[G]danse,
Allez, venez, laissez faire l'insou\[Am]ciance.
\endchorus
\beginverse
\[Am]À deux, à mille je sais qu'on est ca\[F]pable,
Tout est possible, tout est réali\[C]sable.
On peut s'enfuir bien plus haut que nos \[G]rêves,
On peut partir bien plus loin que la \[Am]grève.
\endverse
\beginverse
\[Am]Avec l'envie, la force et le cou\[F]rage,
Le froid, la peur ne sont que des mi\[C]rages.
Laissez tomber les malheurs pour une \[G]fois,
Allez, venez, reprenez avec \[Am]moi.
\endverse
\beginverse
\[Am]Je sais, c'est vrai, ma chanson est na\[F]ïve,
Même un peu bête, mais bien inoffen\[C]sive.
Et même si elle ne change pas le \[G]monde,
Elle vous invite à entrer dans la \[Am]ronde.
\endverse
\beginverse
\[Am]L'espoir, l'ardeur sont tout ce qu'il te \[F]faut,
Mes bras, mon c{\oe}ur, mes épaules et mon \[C]dos.
Je veux te voir des étoiles dans les \[G]yeux,
Je veux nous voir insoumis et heu\[Am]reux.
\endverse
\endsong

\beginsong{Tout le bonheur du monde}[by={Sinsemilia}]
\beginchorus
\[Am]On vous souhaite tout le bonheur du \[F]monde
Et que \[G]quelqu'un vous tende la \[Am]main.
Que vot' \[Am]chemin évite les \[F]bombes,
Qu'il \[G]mène vers de calmes jar\[Am]dins.
\[Am]On vous souhaite tout le bonheur du \[F]monde
Pour au\[G]jourd'hui comme pour de\[Am]main.
Que vot' \[Am]soleil éclaircisse l'\[F]ombre,
Qu'il \[G]brille d'amour au quoti\[Am]dien.
\endchorus
\beginverse
\[G]Puisque l'av'nir vous appar\[Am]tient,
Puis\[Bb]qu'on n'contrôle pas votre destin,
Que votre en\[Dm]vol est pour demain.
\[G]Comme tout c'qu'on a à vous of\[Am]frir
Ne sau\[Bb]rait toujours vous suffire,
Dans cette li\[Dm]berté à venir.
\[Bb]Puisqu'on n'sera pas toujours \[Dm]là,
\[Bb]Comme on le fut aux premiers \[Dm]pas.
\endverse
\beginverse
\[G]Toute une vie s'offre devant \[Am]vous,
Tant \[Bb]d'rêves à vivre jusqu'au bout,
Sûr'ment tant \[Dm]d'joies au rendez-vous.
\[G]Libres de faire vos propres \[Am]choix,
De choi\[Bb]sir quelle sera vot' voie
Et où celle-\[Dm]ci vous emmèn'ra.
\[Bb]J'espère juste que vous prendrez l'\[Dm]temps
\[Bb]De profiter de chaque ins\[Dm]tant.
\endverse
\beginverse
\[G]J'sais pas quel monde on vous laiss'\[Am]ra,
On fait \[Bb]d'not' mieux, seulement parfois,
J'ose espé\[Dm]rer qu'c'la suffira.
\[G]Pas à sauver votre insou\[Am]ciance,
Mais à a\[Bb]paiser not' conscience,
Pour l'reste, j'me \[Dm]dois d'vous faire confiance.
\endverse
\endsong

\beginsong{Un ami}
\beginverse
\[G]Un ami, \[C]c'est ce\[D7]lui qui de\[G]vine
Quand tu \[Em]fais sem\[Am]blant de \[D7]rire
\[G]Un ami, \[C]c'est ce\[D7]lui qui te \[G]tend
La \[Em]main quand tu \[Am]vas tom\[D7]ber
\endverse
\beginchorus
\[C]C'est celui qui \[G]sait ce que tu \[Em]vaux
\[C]C'est celui qui \[G]voit ce que tu \[D7]vaux
\[C]Un a\[D7]mi, \[G]un vrai, ça n'a \[Em]pas de prix
\[C]Un a\[D7]mi, c'est \[G]tout
\endchorus
\endsong

\beginsong{Une belle histoire}[by={Michel Fugain}]
\beginverse
\[Am]C'est un beau ro\[Dm7]man, c'est une belle his\[G]toire, \[Cmaj7]
\[Fmaj7]C'est une romance d'aujour\[Esus4]d'hui. \[E]
\[Am]Il rentrait chez \[Dm7]lui, là-haut vers le \[G]brouillard, \[Cmaj7]
\[Fmaj7]Elle descendait dans le mi\[Esus4]di, le mi\[E]di.
\endverse
\beginchorus
\[Am]Ils se sont trou\[Dm7]vés au bord du che\[G7]min, \[Cmaj7]
\[Fmaj7]Sur l'autoroute des vacances, c'était sans doute un jour de chance.
\[Am]Ils avaient le \[Dm7]ciel à portée de \[G7]main, \[Cmaj7]
\[Fmaj7]Un cadeau de la providence,
Alors pourquoi penser au lende\[Em]main. \[Dm] \[G]
\endchorus
\beginverse
\[Am]Ils se sont ca\[Dm7]chés dans un grand champ de \[G7]blé, \[Cmaj7]
\[Fmaj7]Se laissant porter par le cou\[Esus4]rant. \[E]
\[Am]Se sont racon\[Dm7]té leur vie qui commen\[G7]çait, \[Cmaj7]
\[Fmaj7]Ils n'étaient encore que des en\[Esus4]fants, des en\[E]fants.
\endverse
\beginverse
\[Am]Ils se sont quit\[Dm7]tés au bord du ma\[G]tin, \[Cmaj7]
\[Fmaj7]Sur l'autoroute des vacances, c'était fini le jour de chance.
\[Am]Ils reprirent alors \[Dm7]chacun leur che\[G]min, \[Cmaj7]
\[Fmaj7]Saluèrent la providence
En se faisant un signe de la \[Em]main. \[Dm] \[G]
\endverse
\beginverse
\[Am]Il rentra chez \[Dm7]lui, là-haut vers le \[G]brouillard, \[Cmaj7]
\[Fmaj7]Elle est descendue, là-bas dans le \[Esus4]mi\[E]di. \[Am]
\[Am]C'est un beau ro\[Dm7]man, c'est une belle his\[G]toire, \[Cmaj7]
\[Fmaj7]C'est une romance d'aujour\[Esus4]d'hui. \[E] \[A]
\endverse
\endsong

\beginsong{Un régiment de fromages blancs}
\beginverse
\[G]Un régiment de fro\[D7]mages blancs
Dé\[G]clare la guerre au camem\[D7]bert
Mais le port-sa\[G]lut n'a pas vou\[D7]lu
Car le roque\[G]fort était trop \[D7]fort !
\endverse
\beginverse
\[G]Les marches cre\[D7]vées font les bles\[G]sées
Les asti\[D7]cots hissent le dra\[G]peau
Et ma chan\[D7]son est termi\[G]née
Je vais vous la re\[D7]commen\[G]cer !
\endverse
\endsong

\beginsong{Vent frais}
\textnote{Capo 2}
\beginverse
\[Em]Vent frais, \[D]vent du ma\[Em]tin, \[B7]
\[Em]Vent qui \[D]souffle au sommet des \[G]grands pins. \[B7]
\[Em]Joie du \[D]vent qui \[C]souffle, allons dans le \[B7]grand...
\endverse
\beginverse
\[Em]Vent chaud, \[D]vent du mi\[Em]di, \[B7]
\[Em]Vent qui \[D]souffle au sommet des \[G]épis. \[B7]
\[Em]Joie du \[D]vent qui \[C]souffle, allons dans le \[B7]grand...
\endverse
\beginverse
\[Em]Vent froid, \[D]vent de la \[Em]pluie, \[B7]
\[Em]Vent qui \[D]souffle jusqu'au bout de la \[G]nuit. \[B7]
\[Em]Joie du \[D]vent qui \[C]souffle, allons dans le grand \[B7]vent. \[Em]
\endverse
\endsong

\beginsong{On ira}[by={Zaz}]
\textnote{Capo 1}
\beginverse
\[Am]On ira écouter Harlem au coin de Manhattan,
\[C]On ira rougir le thé dans les souks à Amman,
\[Dm]On ira nager dans le lit du fleuve Sénégal,
\[F]Et on verra brûler Bombay sous un feu de Ben\[G]gale.
\endverse
\beginverse
\[Am]On ira gratter le ciel en dessous de Kyoto,
\[C]On ira sentir Rio battre au c{\oe}ur de Janeiro,
\[Dm]On lèvera nos yeux sur le plafond de la chapelle Sixtine,
\[F]Et on lèvera nos verres dans le café Pouch\[G]kine.
\endverse
\beginverse*
\[Am]Oh qu'elle est belle notre chance
\[C]Aux mille couleurs de l'être humain,
\[Dm]Mélangées de nos différences,
\[F]À la croisée des des\[E]tins.
\endverse*
\beginchorus
\[Am]Vous êtes les étoiles, nous sommes l'univers,
\[C]Vous êtes un grain de sable, nous sommes le désert,
\[Dm]Vous êtes mille pages et moi je suis la plume,
\[F]Oh oh oh oh, oh oh \[G]oh.
\[Am]Vous êtes l'horizon et nous sommes la mer,
\[C]Vous êtes les saisons et nous sommes la terre,
\[Dm]Vous êtes le rivage et moi je suis l'écume,
\[F]Oh oh oh oh, oh oh \[G]oh.
\endchorus
\beginverse
\[Am]On dira que les rencontres font les plus beaux voyages,
\[C]On verra qu'on ne mérite que ce qui se partage,
\[Dm]On entendra chanter des musiques d'ailleurs,
\[F]Et l'on saura donner ce que l'on a de meil\[G]leur.
\endverse
\endsong

\beginsong{On ira}[by={Jean-Jacques Goldman}]
\textnote{Capo 3}
\beginverse
\[G]On partira de nuit, l'heure où l'on \[D]doute
Que demain revienne en\[C]core.
Loin des villes soumises on suivra l'auto\[G]route,
Ensuite on perdra tous les \[D]nords.
\[G]On laiss'ra nos clés, nos cartes et nos \[D]codes,
Prisons pour nous rete\[C]nir.
Tous ces gens qu'on voit vivre comme s'ils igno\[G]raient
Qu'un jour il faudra mou\[D]rir,
Et \[C]qui se font sur\[D]prendre au soir.
\endverse
\beginchorus
Oh \[G]belle, on i\[D]ra,
On par\[Em]tira toi et moi, \[C]où ? je ne sais pas.
\[G]Y a que les routes qui sont \[D]belles,
Et \[C]peu importe où \[D]elles nous mènent.
Oh \[G]belle, on i\[D]ra,
On sui\[Em]vra les étoiles et les cher\[C]cheurs d'or.
\[G]Si on trouve, \[D]on cherchera en\[C]co\[D]re.
\endchorus
\beginverse
\[G]On échappe à rien, pas même à ses \[D]fuites,
Quand on se pose on est \[C]mort.
Oh j'ai tant obéi, si peu choisi pe\[G]tite,
Et le temps perdu me dé\[D]vore.
\[G]On prendra les froids, les brûlures en \[D]face,
On interdira les tiè\[C]deurs,
Les fumées, les alcools et les calmants cui\[G]rasses
Qui nous ont volé nos dou\[D]leurs,
La \[C]vérité nous f'ra plus \[D]peur.
\endverse
\endsong

\beginsong{Contre vents et marées}[by={Caravelles de Vincennes}]
\textnote{Capo 2}
\beginchorus
\[E] \[Am]Contre vents et marées, oser \[F]prendre le large,
\[G]Oser tourner la page, vivre est à inven\[Em]ter.
\[C]Contre vents et marées, se frotter \[Cmaj7]aux orages,
\[Dm]La vie comme un voyage où risquer c'est ai\[E]mer.
\endchorus
\beginverse
\[Am]Même si le ciel est \[E]triste,
Qu't'es \[Dm]seul sur la piste,
Ou l'dernier d'la \[E]liste, la la la.
\[Am]Même si t'attends tou\[E]jours,
Que ce \[Dm]soit ton tour,
Si tu \[E]cries « Au s'cours ! »
\endverse
\beginverse*
\[C]N'oublie pas, n'ou\[G]blie pas,
\[Am]On est tous ensemble sur le même raf\[E]fiot.
\[C]N'oublie pas, n'ou\[G]blie pas,
\[Am]C'est le même vent qui secoue ce ba\[E]teau.
Le \[F]vent, le vent, le vent, le \[G]vent de la vie, \[F]
Le vent, le vent, le \[G]vent de l'Es\[C]prit.
\[C]Viens, viens \[C]changer le pay\[G]sage,
\[Am]Viens, viens repousser les \[E]nuages,
De tes \[F]yeux, de tes yeux,
Peut \[Dm]naître le ciel \[E]bleu.
\endverse*
\beginverse
\[Am]Pour ceux qui sont sans \[E]voix,
Qui \[Dm]n'existent pas,
Dont on ne \[E]veut pas, pa la da.
\[Am]Pour ceux qui sont ban\[E]nis,
Chas\[Dm]sés dans la nuit,
Qui n'ont \[E]que leur cri.
\endverse
\beginverse
\[Am]Même si t'es fati\[E]gué,
Que t'as \[Dm]mal aux pieds,
Que t'en as as\[E]sez, la la la.
\[Am]Même si, découra\[E]gé,
Tu veux \[Dm]t'arrêter,
Tu n'veux \[E]plus marcher.
\endverse
\endsong

\beginsong{L'aventurier}[by={Indochine}]
\beginverse
\[G#m]Égaré dans la vallée infer\[E]nale,
Le \[B]héros s'appelle Bob Mo\[C#]rane.
\[G#m]À la recherche de l'Ombre \[E]Jaune,
Le \[B]bandit s'appelle Mister Kali \[C#]Jones.
\[G#m]Avec l'ami Bill Ballan\[E]tine,
\[B]Sauvé de justesse des croco\[C#]diles.
\[G#m]Stop au trafic des Cara\[E]ïbes,
\[B]Escale dans l'opération Nada\[C#]wieb.
\endverse
\beginchorus
\[Dm]Et soudain surgit face au \[Am]vent
\[C]Le vrai héros de tous les \[G]temps.
\[Dm]Bob Morane contre tout cha\[Am]cal,
\[C]L'aventurier contre tout guer\[G]rier.
\endchorus
\beginverse
\[G#m]Le c{\oe}ur tendre dans le lit de Miss \[E]Clarck,
\[B]Prisonnière du Sultan de Jara\[C#]wak.
\[G#m]En pleine terreur à Manicou\[E]agan,
\[B]Isolé dans la jungle bir\[C#]mane.
\[G#m]Emprisonnant les fli\[E]bustiers,
\[B]L'ennemi est démas\[C#]qué.
\[G#m]On a volé le collier de \[E]Civa,
Le \[B]Maharadjah en répon\[C#]dra.
\endverse
\beginverse
\[G#m]Dérivant à bord du Sam\[E]pang,
\[B]L'aventure au parfum d'Yla\[C#]lang.
\[G#m]Son surnom, Samouraï du So\[E]leil,
En \[B]démantelant le gang de l'Archi\[C#]pel.
\[G#m]L'otage des guerriers du Doc Xha\[E]tan,
Il \[B]s'en sortira toujours à \[C#]temps.
\[G#m]Tel l'aventurier soli\[E]taire,
\[B]Bob Morane est le roi de la \[C#]terre.
\endverse
\endsong

\beginsong{J'ai une tante au Maroc}
\beginchorus
\[C]J'ai une tante qui est au Maroc et qui s'appelle Hip \[G]Hop ! (bis)
\[C]J'ai une tante qui est au \[F]Maroc (bis)
\[C]J'ai une tante qui est au Ma\[G]roc et qui s'appelle Hip \[C]Hop !
\endchorus
\beginverse
\[C]Elle traverse le désert en dromadaire... Houla houla... Hip \[G]Hop ! (bis)
\[C]Elle traverse le dé\[F]sert (bis)
\[C]Elle traverse le désert en droma\[G]daire... Houla houla... Hip \[C]Hop !
\endverse
\beginverse
\[C]Et elle boit d'la grenadine quand elle a soif... Glou glou... Houla houla... Hip \[G]Hop !
\endverse
\beginverse
\[C]Et elle mange du chocolat quand elle a faim... Miam miam... Glou glou... Houla houla... Hip \[G]Hop !
\endverse
\beginverse
\[C]Et elle met une mini-jupe quand elle a chaud... Mini Mini... Miam miam... Glou glou... Houla houla... Hip \[G]Hop !
\endverse
\beginverse
\[C]Et elle tire au pistolet quand elle a peur... Pan Pan... Mini Mini... Miam miam... Glou glou... Houla houla... Hip \[G]Hop !
\endverse
\endsong

\beginsong{Derrière chez ma tante il y a un étang}
\beginverse
\[D]Derrière chez ma tante, il y a un é\[Em]tang, \[A]
\[D]Derrière chez ma tante, il y a un é\[Em]tang. \[A]
\[G]Je me ferai an\[D]guille, anguille dans l'é\[A]tang,
\[G]Je me ferai an\[D]guille, anguille dans l'é\[A]tang. \[D]
\endverse
\beginverse
\[D]Si tu te fais anguille, anguille dans l'é\[Em]tang, \[A]
\[D]Si tu te fais anguille, anguille dans l'é\[Em]tang. \[A]
\[G]Je me ferai pê\[D]cheur, je t'aurai en pê\[A]chant,
\[G]Je me ferai pê\[D]cheur, je t'aurai en pê\[A]chant. \[D]
\endverse
\beginverse
\[D]Si tu te fais pêcheur pour m'avoir en pê\[Em]chant, \[A]
\[G]Je me ferai a\[D]louette, alouette dans les \[A]champs.
\endverse
\beginverse
\[D]Si tu te fais alouette, alouette dans les \[Em]champs, \[A]
\[G]Je me ferai chas\[D]seur, je t'aurai en chas\[A]sant.
\endverse
\beginverse
\[D]Si tu te fais chasseur pour m'avoir en chas\[Em]sant, \[A]
\[G]Je me ferai non\[D]nette, nonnette dans un cou\[A]vent.
\endverse
\beginverse
\[D]Si tu te fais nonnette, nonnette dans un cou\[Em]vent, \[A]
\[G]Je me ferai prê\[D]cheur, je t'aurai en prê\[A]chant.
\endverse
\beginverse
\[D]Si tu te fais prêcheur pour m'avoir en prê\[Em]chant, \[A]
\[G]Je me donn'rai à \[D]toi puisque tu m'aimes \[A]tant. \[D]
\endverse
\endsong

\beginsong{La tribu de Dana}[by={Manau}]
\textnote{Capo 3}
\beginverse
\[Am]Le vent souffle sur les plaines de la Bretagne armo\[G]ricaine,
\[F]Je jette un dernier regard sur ma femme, mon fils et mon do\[E]maine.
\[Am]Akim, le fils du forgeron est venu me cher\[G]cher,
\[F]Les druides ont décidé de mener le combat dans la val\[E]lée. \[Am]
\[Am]C'est l'heure maintenant de défendre notre \[G]terre
\[F]Contre une armée de Cimmériens prête à croiser le \[E]fer.
\[Am]Car c'est la première fois pour moi que je pars au com\[G]bat
\[F]Et j'espère être digne de la tribu de Da\[E]na. \[Am]
\endverse
\beginchorus
\[Am]Dans la vallée oh oh \[G]de Dana lali\[F]lala. \[E]
Dans la vallée oh \[Am]oh j'ai pu entendre les \[G]échos.
\[Am]Dans la vallée oh oh \[G]de Dana lali\[F]lala. \[E]
Dans la vallée oh \[Am]oh des chants de guerre près des tom\[G]beaux. \[Am]
\endchorus
\beginverse
\[Am]Après quelques incantations de druides et de ma\[G]gie,
\[F]Toute la tribu, le glaive en main courait vers l'enne\[E]mi.
\[Am]Mes frères tombaient l'un après l'autre devant mon re\[G]gard,
\[F]Sous le poids des armes que possédaient tous ces bar\[E]bares.
\[Am]Fallait défendre la terre de nos ancêtres enterrés \[G]là
\[F]Et pour toutes les lois de la tribu de Da\[E]na. \[Am]
\endverse
\beginverse
\[Am]Au bout de la vallée on entendait le son d'une \[G]corne,
\[F]D'un chef ennemi qui rappelait toute sa \[E]horde.
\[Am]Le vent souffle toujours sur la Bretagne armori\[G]caine
\[F]Et j'ai rejoint ma femme, mon fils et mon do\[E]maine.
\[Am]J'ai tout reconstruit de mes mains pour en arriver \[G]là,
\[F]Je suis devenu roi de la tribu de Da\[E]na. \[Am]
\endverse
\endsong

\beginsong{Les lacs du Connemara}[by={Michel Sardou}]
\beginchorus
\[D]Terre brûlée au vent des landes de \[D]pierre,
Autour des lacs, c'est pour les vi\[Bm]vants
Un peu d'enfer, le Conne\[A]mara.
\[Em]Des nuages noirs qui \[B7]viennent du \[Em]nord
Colorent la \[B7]terre, les lacs les ri\[C]vières :
C'est le décor du Conne\[A]mara.
\endchorus
\beginverse
\[D]Au printemps suivant, le \[G]ciel irlan\[D]dais
Était en paix, Mau\[A]reen a plongé \[Bm]nue
Dans un lac du Conne\[Asus4]ma\[A]ra. \[B7]
\[Em]Sean Kelly s'est dit : « Je \[B7]suis catho\[Em]lique.
Maureen aus\[B7]si », l'église en gra\[Em]nit
De \[B7]Lime\[Em]rick, Maureen a dit \[Em7]« oui ». \[A]
\endverse
\beginverse
\[F]De Tiperrary, Bally-\[Bb]Connelly et de \[F]Galway,
Ils \[C]sont arrivés dans le \[Dm]comté du Conne\[Csus4]ma\[C]ra.
\[Gm]Y'avait les Connor, les O'\[D7]Conolly, les Fla\[Gm]herty
Du \[D7]Ring of Kerry et de \[Gm]quoi boire trois jours et deux \[C]nuits. \[Ab7]
\endverse
\beginverse*
Là-\[Db]bas, au Conne\[Ab]ma\[Db]ra, \[Bb7]
On \[Ebm7]sait tout le prix du si\[Ebsus4]len\[Db]ce. \[Ab7]
Là-\[F7]bas, \[Bbm]au Conne\[F7]mara, on \[Gb]dit que la vie
C'est une fo\[F]lie et que la fo\[Ebm7]lie, ça se \[Ab]danse.
\endverse*
\beginverse
\[D]On y vit encore au \[G]temps des \[D]Gaels
Et de Crom\[A]well, au rythme des \[Bm]pluies
Et du soleil, au pas des che\[Asus4]vaux. \[A]
\[Em]On y croit encore aux \[B7]monstres des \[Em]lacs
Qu'on voit na\[B7]ger certains soirs d'é\[Em]té
Et replonger pour l'éter\[A]nité.
\endverse
\endsong

\beginsong{À nos souvenirs}[by={Trois Cafés Gourmands}]
\beginverse
\[C#m]Comment puis-je oublier ce \[A]coin de paradis ?
Ce \[E]petit bout de terre où \[B]vit encore mon père.
\[C#m]Comment pourrais-je faire pour \[A]me séparer d'elle ?
Ou\[E]blier qu'on est frères, \[B]belle Corrèze charnelle.
\endverse
\beginverse
\[C#m]Oublier ce matin que \[A]tu es parisien,
Que t'as \[E]de l'eau dans l'vin, que \[B]tu es parti loin ?
\[C#m]Ce n'était pas ma faute, \[A]on joue des fausses notes,
On se \[E]trompe de chemin et \[B]on a du chagrin.
\endverse
\beginverse
On se joue tout un \[C#m]drame, on a des vagues à \[A]l'âme
Tu as du mal au \[E]coeur, tu as peur du bon\[B]heur
Acheter des ta\[C#m]bleaux et des vaches en pho\[A]to
C'est tout c'que t'as trou\[E]vé, pour te la rappe\[B]ler
\endverse
\beginverse
Vous me trouvez un peu \[C#m]con, n'aimez pas ma chan\[A]son
Vous me croyez bi\[E]zarre, un peu patrio\[B]tard
Le fruit de ma réfle\[C#m]xion ne touchera per\[A]sonne
Si vos pas ne ré\[E]sonnent jamais dans ma ré\[B]gion
\endverse
\beginverse
\[C#m]C'est pire qu'une religion, \[A]au-delà d'une confession,
Je \[E]l'aime à en mourir, \[B]pour le meilleur et pour le pire.
Et si je monte au \[C#m]ciel il y aura peut être \[A]Joel
Guillaume et Jéré\[E]my et mon cousin Pie\[B]dri
\endverse
\beginverse
Yoan s'ra en vo\[C#m]yage dans un autre \[A]pays
Allez fais tes ba\[E]gagesm viens rejoindre tes a\[B]mis
On veut du Claudie mu\[C#m]sette à en perdre la \[A]tête
On veut un dernier cha\[E]brol, un petit coup de \[B]niole 
\endverse
\beginverse
\[C#m]Les yeux de nos grands-mères, \[A]la voix de nos grands-pères,
L'o\[E]deur de cette terre, \[B]vue sur les Monédières.
C'est pire qu'un testa\[C#m]ment, au delà d'une confi\[A]dence
On est des p'tits en\[E]fants de ce joli coin de \[B]France
\endverse
\beginverse
Enterrez nous vi\[C#m]vants, baillonnés s'il le \[A]faut
Mais prenez soin a\[E]vant de remplir not' ja\[B]bot
\endverse
\beginverse 
La relève est pour \[C#m]toi, notre petit Lu\[A]cas
On t'laisse en héri\[E]tage la piste nous on dé\[B]gage
Le temps nous a gâ\[C#m]té, on en a bien profi[A]té
On a des souvenirs en \[E]tête, ce soir, faisons la \[B]fête !
\endverse
\beginchorus
Acceptez ma ren\[C#m]gaine, elle veut juste dire « je t'\[A]aime »,
Soyez sûrs, j'en suis \[E]fier, j'ai la Corrèze en cathé\[B]ter.
D'être avec vous ce \[C#m]soir, j'ai le c{\oe}ur qui pé\[A]tille,
Mimi sers-nous à \[E]boire, on a les yeux qui \[B]brillent...
\endchorus
\beginverse
Papaya paya pa\[C#m]pa paya paya pa\[A]pa   | x3
Papaya paya pa\[E]pa paya paya pa\[B]pa     |
\endverse
\beginchorus
Acceptez ma ren\[C#m]gaine, elle veut juste dire « je t'\[A]aime »,
Soyez sûrs, j'en suis \[E]fier, j'ai la Corrèze en cathé\[B]ter.
D'être avec vous ce \[C#m]soir, j'ai le c{\oe}ur qui pé\[A]tille,
Mimi sers-nous à \[E]boire, on a les yeux qui \[B]brillent...
\endchorus
\endsong

\beginsong{Petrouchka}[by={Mannick}]
\textnote{Capo 3}
\beginverse
\[Am]Petrouchka, ne pleure pas,
\[Dm]Entre vite dans la \[Am]ronde,
\[Dm]Fais danser tes nattes \[Am]blondes,
\[E7]Ton petit chat revien\[Am]dra.
\[Dm]Il s'est fait polichi\[Am]nelle
\[Dm]Dans les chemises en den\[Am]telle
De \[E7]ton grand-pa\[Am]pa.
\endverse
\beginchorus
\[F]Tant que chante la co\[G]lombe
Par-\[C]dessus \[E7]le \[Am]toit,
\[Dm]Danse avant que la nuit \[E7]tombe,
Jolie \[Am]Pe\[G]trouch\[C]ka.
\[F]Tant que chante la co\[G]lombe
Par-\[C]dessus \[E7]le \[Am]toit,
\[Dm]Danse avant que la nuit \[E7]tombe,
Jolie \[F]Pe\[E7]trouch\[Am]ka.
\endchorus
\beginverse
\[Am]Petrouchka, ne pleure pas,
\[Dm]Mets ton grand fichu de \[Am]laine,
\[Dm]Viens avec nous dans la \[Am]plaine,
\[E7]Ton petit chat revien\[Am]dra.
\[Dm]Il fait quatre gali\[Am]pettes,
Se \[Dm]déguise en marion\[Am]nette
Dès \[E7]que tu t'en \[Am]vas.
\endverse
\beginverse
\[Am]Petrouchka, ne pleure pas,
\[Dm]Puisqu'il aime la mu\[Am]sique,
\[Dm]Chante-lui cet air ma\[Am]gique,
\[E7]Ton petit chat revien\[Am]dra.
\[Dm]Il nous dansera peut-\[Am]être
Sur le \[Dm]bord de la fe\[Am]nêtre
Une \[E7]mazur\[Am]ka.
\endverse
\endsong

\beginsong{J'veux du soleil}[by={Au P'tit Bonheur}]
\beginverse
\[Bm]J'suis resté qu'un enfant
Qu'au\[F#7]rait grandi trop vite,
Dans un \[Em]monde en super plastique,
\[F#7]J'veux retrouver maman.
\[Bm]Qu'elle m'raconte des histoires
De \[F#7]Jane et de Tarzan,
De \[Em]princesses et de cerf-volants,
\[F#7]J'veux du soleil dans ma mémoire.
\endverse
\beginchorus
\[Bm]J'veux du so\[F#7]leil,
J'veux du so\[Em]leil,
J'veux du so\[F#7]leil !
\endchorus
\beginverse
\[Bm]J'veux traverser les océans,
De\[F#7]venir Monte-Cristo,
Au \[Em]clair de lune m'échapper
\[F#7]D'la citadelle.
\[Bm]J'veux devenir roi des marécages,
\[F#7]M'sortir de ma cage,
Un \[Em]père Noël pour Cendrillon
\[F#7]Sans escarpin.
\endverse
\beginverse
\[Bm]J'veux faire danser maman
Au \[F#7]son clair des grillons,
J'veux \[Em]retrouver mon sourire d'enfant
\[F#7]Perdu dans l'tourbillon.
\[Bm]Dans le tourbillon de la vie
Qui \[F#7]fait que l'on oublie
Qu'\[Em]l'on est restés des mômes
\[F#7]Bien au fond de nos abris.
\endverse
\endsong

\beginsong{Sur ma route}[by={Black M}]
\textnote{Capo 3}
\beginchorus
\[Em]Sur ma route, oui, il y a eu du \[C]move, oui,
De l'aventure dans l'\[G]movie, une vie de \[D]roots.
\[Em]Sur ma route, oui, je n'compte plus les \[C]soucis,
De quoi devenir \[G]fou, oui, une vie de \[D]roots.
\endchorus
\beginverse
\[Em]Sur ma route, j'ai eu des moments de \[C]doute,
J'marchais sans savoir vers \[G]où, j'étais têtu, rien à \[D]foutre.
\[Em]Sur ma route, j'avais pas de bagages en \[C]soute,
Et dans ma poche pas un \[G]sou, juste la famille entre \[D]nous.
\endverse
\beginverse
\[Em]Sur ma route, on m'a fait des coups en \[C]douce,
L'impression qu'mon c{\oe}ur en \[G]souffre, mais je suis sous anes\[D]thésie.
\[Em]Sur mon ch'min, j'ai croisé pas mal d'an\[C]ciens,
Ils me parlaient du lende\[G]main et que tout allait si \[D]vite.
\endverse
\endsong

\beginsong{La Bohème}[by={Charles Aznavour}]
\beginverse
\[Fm]Je vous parle d'un temps que les moins de vingt ans
Ne peuvent pas con\[Cm]naître. Montmartre en ce temps-\[Fm]là
Accrochait ses lilas jusque sous nos fe\[Cm]nêtres.
Et si l'humble gar\[Fm]ni qui nous servait de nid
Ne payait pas de \[Cm]mine, c'est là qu'on s'est con\[Fm]nus,
Moi qui criait fa\[G9]mine et toi qui posais \[Cm]nue.
\endverse
\beginchorus
\[C7]La bo\[Fm]hème, la bo\[Cm]hème,
Ça vou\[Fm]lait dire on est heu\[G7]reux. \[Cm] \[C7]
La bo\[Fm]hème, la bo\[Cm]hème,
Nous ne man\[Fm]gions qu'un jour sur \[G7]deux. \[Cm]
\endchorus
\beginverse
\[Fm]Dans les cafés voisins nous étions quelques-uns
Qui attendions la \[Cm]gloire, et bien que misé\[Fm]reux
Avec le ventre creux nous ne cessions d'y \[Cm]croire.
Et quand quelque bis\[Fm]tro contre un bon repas chaud
Nous prenait une \[Cm]toile, nous récitions des \[Fm]vers
Groupés autour du \[G9]poêle en oubliant l'hi\[Cm]ver.
\endverse
\beginverse
\[Fm]Souvent il m'arrivait devant mon chevalet
De passer des nuits \[Cm]blanches, retouchant le des\[Fm]sin
De la ligne d'un sein, du galbe d'une \[Cm]hanche.
Et ce n'est qu'au ma\[Fm]tin qu'on s'asseyait enfin
Devant un café-\[Cm]crème, épuisés mais ra\[Fm]vis,
Fallait-il que l'on \[G9]s'aime et qu'on aime la \[Cm]vie.
\endverse
\beginverse
\[Fm]Quand au hasard des jours je m'en vais faire un tour
À mon ancienne a\[Cm]dresse, je ne reconnais \[Fm]plus
Ni les murs ni les rues qui ont vu ma jeu\[Cm]nesse.
En haut d'un esca\[Fm]lier je cherche l'atelier
Dont plus rien ne sub\[Cm]siste, dans son nouveau dé\[Fm]cor
Montmartre semble \[G9]triste et les lilas sont \[Cm]morts.
\endverse
\endsong

\beginsong{La route}[by={Jamboree}]
\beginverse
\[Am]Je me rappelle ces jours de liesse
Où l'on partait vers l'inno\[Em]cence,
Sur tes \[Bm]chemins, le pas pressé et \[Em]fier.
\[Am]Je me rappelle ces longs chemins
Dont on ne voyait pas la \[Em]fin,
Ces \[Bm]douces rencontres au goût de miel,
Ce \[Em]vent si pur, cette mer si \[D]belle.
\endverse
\beginverse
\[G]Je me rappelle de nos chansons,
De ces \[D]décors aux vastes horizons,
De nos sou\[Em]rires, de nos douleurs,
De nos sou\[Bm]pirs et de nos \[C]peurs. \[D]
\endverse
\beginchorus
\[G]Vagabondes, nos âmes ont pris la \[D]route,
Adieu c{\oe}ur dessé\[Em]ché, on trouv'ra bien d'quoi t'arro\[Bm]ser.
Adieu co\[C]lère, adieu désirs a\[G]mers,
On a op\[D]té pour la folie passagère.
\endchorus
\beginverse
\[Am]Nouveaux rivages sans grillages,
\[Em]Nouvelles terres sans barrières.
\[Am]La route nous mène où le vent nous en\[Em]traîne,
\[Bm]Prenons-nous la main et oublions de\[Em]main,
La \[D]route nous guidera, on verra bien là-bas.
\endverse
\endsong

\beginsong{Y'a des zazous}[by={Andrex}]
\beginverse
\[F]Jusqu'ici sur \[F/C]terre un \[F]homme pouvait \[F/C]être
\[F]Blanc ou noir ou \[F/C]rouge, ou jaune et puis c'est \[C7]tout.
Mais une \[Gm]autre race est en \[C7]train d'appa\[Gm]raître,
C'est les \[C]Zazous, \[C7]c'est les Za\[F]zous !
\endverse
\beginchorus
\[C7]Y a des Zazous dans mon \[F/C]quartier,
\[F]Moi je l'suis déjà à moi\[C7]tié.
À votre \[Gm]tour un de ces \[C7]jours
Vous serez \[G5]tous Zazous comme eux,
Car le \[G7]Zazou c'est conta\[C]gieux. \[C7]
Ça commence par un trem\[F/C]blement
\[F]Qui vous prend soudain brusque\[C7]ment,
Et puis on \[Gm]pousse des hurle\[F]ments ! \[F/C]
\[C7]Wa da la di zou, za zi zou \[F]la !
\endchorus
\beginverse
\[F]Si vous rencon\[F/C]trez un \[F]jour sur votre pas\[F/C]sage
\[F]Un particu\[F/C]lier coiffé d'un fromage \[C7]mou,
Tenant dans ses \[Gm]doigts un poisson \[C7]dans une \[Gm]cage,
C'est un \[C]Zazou, \[C7]c'est un Za\[F]zou !
\endverse
\endsong

\beginsong{1987}[by={Calogero}]
\beginverse
\[D]Tu te sou\[Dsus4]viens,
Les couleurs sur les bas\[Bm6]kets,
Les \[Bm]crayons dans les cas\[A6]settes,
Je rem\[A]bobine.
\[D]Tu te sou\[Dsus4]viens,
Tous ces rêves plein nos dis\[Bm6]quettes,
À \[Bm]Paris c'était les \[A6]States,
Mille \[A]neuf cent quatre-vingt-\[Em]sept. \[A]
\endverse
\beginchorus
\[D]Il y a certains jours
Où \[A]je reprends mon skate,
\[Bm]Et je vais faire un tour
En \[F#m]mille neuf cent quatre-vingt-\[G]sept.
\[D]Il y a certains jours
Dans \[A]lesquels je me jette,
\[Bm]Et je suis de retour
En \[F#m]mille neuf cent quatre-vingt-\[G]sept.
\[D]Tu sais de tous ces jours
Y'a \[A]rien que je regrette,
\[Bm]Mais parfois je retourne
En \[F#m]mille neuf cent quatre-vingt-\[G]sept.
\endchorus
\beginverse
\[D]Tu te sou\[Dsus4]viens,
Les survêt et les hou\[Bm6]pettes,
Sa\[Bm]brina et sept sur \[A6]sept,
Dans \[A]la cuisine.
\[D]C'était \[Dsus4]rien,
Que douze mois sur la pla\[Bm6]nète,
\[Bm]L'URSS, \[A6]INXS,
On \[A]chantait I want your \[Em]sex. \[A]
\endverse
\endsong

\beginsong{Le chant des sirènes}[by={Fréro Delavega}]
\textnote{Capo 7}
\beginverse
\[Dm]Enfants des parcs, gamins des \[Am]plages,
Le vent menace les châteaux de \[G]sable façonnés de mes \[Dm]doigts.
Le temps n'épargne per\[Am]sonne hélas,
Les années passent, l'écho s'évade \[G]sur la Dune du Pyla.
\endverse
\beginverse*
\[F]Au gré des saisons, des \[Am]photomatons,
Je m'aban\[C]donne à ces lueurs d'autre\[G]fois.
\[F]Au gré des saisons, des \[Am]décisions, \[C]je m'aban\[G]donne.
\endverse*
\beginchorus
\[Dm]Quand les souvenirs s'en \[Am]mêlent, les larmes me \[C]viennent,
Et le chant des sirènes me re\[G]plonge en hiver.
\[Dm]Oh mélancolie cru\[Am]elle, harmonie flu\[C]ette, euphorie soli\[G]taire.
\endchorus
\beginverse
\[Dm]Combien de farces, combien de \[Am]frasques,
Combien de traces, combien de \[G]masques avons-nous laissé là-\[Dm]bas.
Poser les armes, prendre le \[Am]large,
Trouver le calme dans ce va\[G]carme avant que je ne m'y noie.
\endverse
\endsong

\beginsong{Ciel}[by={Maître Gims}]
\beginchorus
\[Am]Ciel, tu m'avais pas dit qu'c'était une magi\[G]cienne, \[Dsus2]
\[Am]Ça n'm'étonnerait pas qu'elle me fasse des \[G]siennes. \[Dsus2]
\[Am]J'ai fait un cauchemar, j'étais à décou\[G]vert, \[Dsus2]
\[Am]J'ai retrouvé la vue dans le Ferrari \[G]vert. \[Dsus2]
\[Am]Elle est tombée du \[G]ciel, \[Dsus2]
\[Am]Oui, c'est une magi\[G]cienne. \[Dsus2]
\endchorus
\beginverse
\[Am]Tu dis que tu l'aimes mais c'est pour de \[G]faux, \[Dsus2]
\[Am]T'as envie de changer le cours des \[G]choses. \[Dsus2]
\[Am]T'as suivi ton âme au fin fond d'la \[G]fosse, \[Dsus2]
\[Am]Aujourd'hui, t'as vu ton plus grand dé\[G]faut. \[Dsus2]
\[Am]C'est-à-dire que la vie n'est pas comme tu \[G]crois, \[Dsus2]
\[Am]Ça veut dire que tout n'est qu'une question de \[G]choix. \[Dsus2]
\[Am]C'est fini mais je garde les plus belles \[G]images. \[Dsus2]
Va leur dire que c'était une belle histoire.
\endverse
\beginverse
\[Am]Tu me suis partout et \[G]même dans mes ga\[Dsus2]lères,
\[Am]J't'ai menti le jour où on \[G]regardait la \[Dsus2]mer.
\[Am]Y a plus rien à faire vu les \[G]bouts de \[Dsus2]verre,
\[Am]Hanté par les doutes et la \[G]peur de te \[Dsus2]perdre.
\endverse
\endsong

\beginsong{Elle me dit}[by={MIKA}]
\textnote{Capo 6}
\beginverse
\[E]Elle me dit, écris une chanson con\[B]tente,
Pas une chanson dépri\[A]mante,
\[E]Une chanson que tout le monde \[B]aime. \[A]
\[E]Elle me dit, tu deviendras million\[B]naire,
Tu auras de quoi être \[A]fier,
\[E]Ne finis pas comme ton \[B]père. \[A]
\endverse
\beginverse
\[E]Elle me dit, ne t'enferme pas dans ta \[B]chambre,
Vas-y secoue-toi et \[A]danse,
\[E]Dis-moi c'est quoi ton pro\[B]blème. \[A]
\[E]Elle me dit, qu'est-ce que t'as, t'as pas l'air \[B]coincé ?
T'es défoncé ou t'es \[A]gay ?
\[E]Tu finiras comme ton \[B]frère. \[A]
\endverse
\beginverse*
\[E]Elle me dit, c'est ta vie,
\[B]Fais c'que tu veux, tant pis.
\[C#m7]Un jour tu comprendras,
\[A]Un jour tu t'en voudras.
\[E]Elle me dit, t'es trop nul,
\[B]Sors un peu de ta bulle.
\[C#m7]Tu fais n'importe quoi,
\[A]On dirait que t'aimes ça.
\endverse*
\beginchorus
\[E]Pourquoi tu gâches ta \[B]vie,
Pourquoi tu gâches ta \[A]vie,
\[E]Pourquoi tu gâches ta vie,
\[B]Danse, danse, \[A]danse !
\endchorus
\beginverse
\[E]Elle me dit, fais comme les autres gar\[B]çons,
Vas taper dans un bal\[A]lon,
\[E]Tu deviendras popu\[B]laire. \[A]
\[E]Elle me dit, un jour je serai plus \[B]là,
Mais c'est quand elle me dit \[A]ça
\[E]Qu'elle me dit un truc que \[B]j'aime. \[A]
\endverse
\beginverse
\[E]Elle me dit, t'as pas encore des ch'veux \[B]blancs
Mais t'auras bientôt 30 \[A]ans,
\[E]Faudrait mieux que tu te ré\[B]veilles. \[A]
\[E]Elle me dit, oui un jour tu me tue\[B]ras,
Mais c'est quand elle me dit \[A]ça
\[E]Qu'elle me dit un truc que \[B]j'aime. \[A]
\endverse
\endsong

\beginsong{Emmenez-moi}[by={Charles Aznavour}]
\beginverse
\[Am]Vers les docks où le poids et l'en\[G]nui
Me courbent le \[Am]dos, \[E7]
\[Am]Ils arrivent le ventre alourdi de \[G]fruits,
Les ba\[Am]teaux. \[E7] \[Am]
\[F]Ils viennent du bout du \[G]monde
Apportant avec \[F]eux des idées vaga\[G]bondes
Aux reflets de ciel \[F]bleu, de mi\[C]rage.
\[C]Traînant un parfum poi\[F]vré de pays incon\[C]nus
Et d'éternels é\[F]tés où l'on vit presque \[C]nu
Sur les \[E7]plages.
\[Am]Moi qui n'ai connu toute ma \[G]vie
Que le ciel du \[Am]Nord, \[E7]
\[Am]J'aimerais débarbouiller ce \[G]gris
En virant de \[Am]bord.
\endverse
\beginchorus
\[Am] \[E7] \[Am] \[G]Emmenez-moi au bout de la \[C]terre,
\[G]Emmenez-moi au pays des mer\[C]veilles,
\[E7]Il me semble que la mi\[Am]sère
\[F]Serait moins \[E7]pénible au so\[Am]leil.
\endchorus
\beginverse
\[Am]Dans les bars à la tombée du \[G]jour
Avec les ma\[Am]rins, \[E7]
\[Am]Quand on parle de filles et d'a\[G]mour,
Un verre à la \[Am]main. \[E7] \[Am]
\[F]Je perds la notion des \[G]choses
Et soudain ma pen\[F]sée m'enlève et me dé\[G]pose
Un merveilleux é\[F]té sur la \[C]grève.
\[C]Où je vois tendant les \[F]bras, l'amour qui comme un \[C]fou
Court au-devant de \[F]moi, et je me pends au \[C]cou
De mon \[E7]rêve.
\[Am]Quand les bars ferment, que les ma\[G]rins
Rejoignent leur \[Am]bord, \[E7]
\[Am]Moi je rêve encor' jusqu'au ma\[G]tin,
Debout sur le \[Am]port.
\endverse
\beginverse
\[Am]Un beau jour, sur un rafiot cra\[G]quant
De la coque au \[Am]pont, \[E7]
\[Am]Pour partir, je travaillerai \[G]dans
La soute à char\[Am]bon. \[E7] \[Am]
\[F]Prenant la route qui \[G]mène
À mes rêves d'en\[F]fant, sur des îles loin\[G]taines,
Où rien n'est impor\[F]tant que de \[C]vivre.
\[C]Où les filles alangui\[F]es vous ravissent le \[C]c{\oe}ur
En tressant, m'a-t-on \[F]dit, de ces colliers de \[C]fleurs
Qui en\[E7]ivrent.
\[Am]Je fuirai, laissant là mon pas\[G]sé,
Sans aucun re\[Am]mords, \[E7]
\[Am]Sans bagage et le c{\oe}ur libé\[G]ré,
En chantant très \[Am]fort.
\endverse
\endsong

\beginsong{Formidable}[by={Stromae}]
\beginchorus
\[Em]Formidable, formi\[D/F#]dable, \[G]
\[C]Tu étais formi\[Am]dable, j'étais fort mi\[Em]nable,
\[B7/F#]Nous étions formi\[Em]dables.
\endchorus
\beginverse
\[Em]Oh bébé, oups, made\[D/F#]moiselle, \[G]
J'vais pas vous dra\[C]guer, promis ju\[Am]ré, j'suis céliba\[Em]taire
Depuis hier pu\[B7/F#]tain, j'peux pas faire d'en\[Em]fant...
\[Em]Cinq minutes quoi ! J't'ai \[D/F#]pas insultée, \[G]
J'suis poli, cour\[C]tois, et un peu fort bour\[Am]ré,
Et pour les mecs comme \[Em]moi, vous avez autre chose à \[B7/F#]faire, hein,
Vous m'auriez vu \[Em]hier...
\endverse
\beginverse
\[Em]Eh, tu t'es regar\[D/F#]dé, tu t'crois \[G]beau
Parce que tu t'es ma\[C]rié, mais c'est qu'un an\[Am]neau.
Mec, t'emballe \[Em]pas, elle va t'larguer comme elles le \[B7/F#]font chaque \[Em]fois.
Et puis l'autre \[Em]fille, tu lui en as par\[D/F#]lé ? \[G]
Si tu veux je lui \[C]dis, comme ça c'est ré\[Am]glé,
Et au p'tit aus\[Em]si, enfin si vous en a\[B7/F#]vez,
Attends trois ans, sept \[Em]ans, et là vous verrez...
\endverse
\beginverse
\[Em]Et petite, oh par\[D/F#]don, petit, \[G]
Tu sais dans la \[C]vie y'a ni méchant ni gen\[Am]til.
Si Maman est \[Em]chiante, c'est qu'elle a peur d'être ma\[B7/F#]mie,
Si Papa trompe Ma\[Em]man, c'est parce que Maman vieillit.
Pourquoi t'es tout \[D/F#]rouge ? Ben reviens ga\[G]min !
Et qu'est-ce que vous avez \[C]tous, à me regarder comme un \[Am]singe, vous ?
Ah oui vous êtes \[Em]saints, vous, bande de ma\[B7/F#]caques !
Donnez-moi un bébé \[Em]singe, il sera...
\endverse
\endsong

\beginsong{L'Amérique pleure}[by={Les Cowboys Fringants}]
\beginverse
\[F]Encore un jour à se lever
En même temps que le so\[C]leil,
La face encore un peu poquée
D'mon quatre heures de sommeil.
\[F]J'tire une coup' de poffes de clope,
Job done pour les vita\[C]mines,
Pis un bon café à l'eau de moppe,
Histoire de s'donner meilleure mine.
\endverse
\beginverse*
\[G] \[E] \[F] \[C] \[G] \[F]
\endverse*
\beginverse
\[F]J'prends le Florida Turnpike,
Pis demain soir j'ta Mont\[C]magny,
Non trucker c'pas vraiment l'Klondike
Mais tu vois du pays.
\[F]Surtout que ça t'fait réaliser
Que derrière les beaux pay\[C]sages,
Y'a tellement d'inégalités
Et de souffrance sur les visages.
\endverse
\beginchorus
\[G]La question qu'j'me pose tout l'\[E]temps,
Mais comment font tous ces \[F]gens ?
Pour croire encore en la \[C]vie,
Dans cette \[G]hypocri\[F]sie.
C'est si triste que des \[G]fois,
Quand je rentre à la mai\[E]son
Et que j'park mon vieux ca\[F]mion,
J'vois toute l'Amérique qui \[C]pleure
Dans mon \[G]rétrovi\[F]seur.
\endchorus
\beginverse
\[F]Moi je traîne dans ma remorque
Tous les excès d'mon é\[C]poque,
La surabondance surgelée,
Shootée, suremballée.
\[F]Pendant qu'les v{\oe}ux pieux passent dans l'beurre,
Que notre insouciance est re\[C]pue,
C'est dans le fond des containers
Que pourront pourrir les surplus.
\endverse
\beginverse*
\[G]La question qu'j'me pose tout l'\[E]temps,
Mais que feront nos en\[F]fants ?
Quand il ne restera \[C]rien,
Que des \[G]ruines et la \[F]faim.
C'est si triste que des \[G]fois,
Quand je rentre à la mai\[E]son
Et que j'parke mon vieux ca\[F]mion,
J'vois toute l'Amérique qui \[C]pleure
Dans mon \[G]rétrovi\[F]seur.
\endverse*
\beginverse
\[F]Sur l'Interstate 95
Partent en fumée tous les \[C]rêves,
Un char en feu dans une bretelle,
Un accident mortel.
\[F]Et au milieu de ce bouchon,
Pas de respect pour la \[C]mort,
Chacun son tour joue du klaxon,
Tellement pressé d'aller nulle part.
\endverse
\beginverse*
\[G]La question qu'j'me pose tout l'\[E]temps,
Mais où s'en vont tous ces \[F]gens ?
Y'a tellement de chars par\[C]tout,
Le monde est \[G]rendu \[F]fou.
C'est si triste que des \[G]fois,
Quand je rentre à la mai\[E]son
Et que j'parke mon vieux ca\[F]mion,
J'vois toute l'Amérique qui \[C]pleure
Dans mon \[G]rétrovi\[F]seur.
\endverse*
\beginverse
\[F]Un aut' truck stop d'autoroute,
Pogné pour manger d'la ch'\[C]noute.
C'est vrai que dans la soupe du jour
Y'a pu tellement d'amour.
\[F]On a tué la chaleur humaine
Avec le service à la \[C]chaîne.
À la télé un aut' malade
Vient d'déclencher une fusillade.
\endverse
\beginverse*
\[G]La question qu'j'me pose tout l'\[E]temps,
Mais comment font ces pauvres \[F]gens
Pour traverser tout le \[C]cours
D'une vie \[G]sans a\[F]mour.
C'est si triste que des \[G]fois,
Quand je rentre à la mai\[E]son
Pis que j'parke mon vieux ca\[F]mion,
Je vois toute l'Amérique qui \[C]pleure
Dans mon \[G]rétrovi\[F]seur.
\endverse*
\beginverse
\[F]Rien n'empêche que moi aussi,
Quand j'roule tout seul dans la \[C]nuit,
J'me demande des fois ce que je fous ici,
Pris dans l'arrière-pays.
\[F]J'pense à tout c'que j'ai manqué
Avec Mimi pi les deux \[C]filles,
Et j'ai ce sentiment fucké
D'être étranger dans ma famille.
\endverse
\beginverse*
\[G]La question que j'me pose tout le \[E]temps,
Pourquoi travailler au\[F]tant,
M'éloigner de ceux que \[C]j'aime,
Tout ça pour \[G]jouer la \[F]game.
C'est si triste que des \[G]fois,
Quand j'suis loin de la mai\[E]son,
Assis dans mon vieux ca\[F]mion,
J'ai toute l'Amérique qui \[C]pleure
Quelque part au \[G]fond du c{\oe}\[F]ur.
\endverse*
\endsong

\beginsong{Les démons de minuit}[by={Images}]
\textnote{Capo 2}
\beginverse
\[G]Rue déserte, dernière cigarette plus rien ne \[Bm]bouge,
\[G]Juste un bar éclaire le comptoir d'un néon \[Bm]rouge.
\[Bb]J'ai besoin de trouver quelqu'un, j'veux pas dor\[Dm]mir,
\[Bb]Je cherche un peu de cha\[Gm]leur à mettre dans mon \[A]c{\oe}ur.
\endverse
\beginchorus
\[Dm]Ils m'entraînent au bout de la \[Bb]nuit, les \[Gm]démons de mi\[C]nuit,
\[Dm]M'entraînent jusqu'à l'insom\[Bb]nie... les \[Gm]fantômes de l'en\[C]nui. \[Dm]
\endchorus
\beginverse
\[G]Dans mon verre je regarde la mer qui se ba\[Bm]lance,
\[G]J'veux un disque de funky music faut que ça \[Bm]danse.
\[Bb]J'aime cette fille sur talons aiguilles qui se ba\[Dm]lance,
\[Bb]Ça met un peu de cha\[Gm]leur au fond de mon \[A]c{\oe}ur.
\endverse
\endsong

\beginsong{Les prisons de Nantes}
\beginverse
\[Am]Dans les prisons de Nantes, \[G]landidididoudan...
\[Am]Dans les prisons de Nantes, y'avait un prisonnier, \[G]y'avait un pri\[Am]sonnier.
\endverse
\beginverse
\[Am]Personne ne le vint l'vouer, \[G]landidididoudan...
\[Am]Personne ne le vint l'vouer, que la fille du geôlier, \[G]la fille du \[Am]geôlier.
\endverse
\beginverse
\[Am]Un jour il lui demande, \[G]landidididoudan...
\[Am]Un jour il lui demande, oui que dit-on de moué, \[G]que dit-on de \[Am]moué.
\endverse
\beginverse
\[Am]On dit de vous en ville, \[G]landidididoudan...
\[Am]On dit de vous en ville, que vous serez pendu, \[G]et vous se\[Am]rez \[Em]pen\[Am]du.
\endverse
\beginverse
\[Am]Mais s'il faut qu'on me pende, \[G]landidididoudan...
\[Am]Mais s'il faut qu'on me pende, déliez-moi les pieds, \[G]déliez-moi \[Am]les \[Em]pieds. \[Am]
\endverse
\beginverse
\[Am]La fille était jeunette, \[G]landidididoudan...
\[Am]La fille était jeunette, les pieds lui a délié, \[G]les pieds lui \[Am]a \[Em]dé\[Am]lié.
\endverse
\beginverse
\[Am]Le prisonnier alerte, \[G]landidididoudan...
\[Am]Le prisonnier alerte, dans la Loire s'est jeté, \[G]dans la Loire \[Am]s'est \[Em]je\[Am]té.
\endverse
\beginverse
\[Am]Je chante pour les belles, \[G]landidididoudan...
\[Am]Je chante pour les belles, surtout celle du geôlier, \[G]celle du \[Am]ge\[Em]ô\[Am]lier.
\endverse
\beginverse
\[Am]Si je reviens à Nantes, \[G]landidididoudan...
\[Am]Si je reviens à Nantes, oui je l'épouserai, \[G]je l'épou\[Am]se\[Em]rai. \[Am]
\endverse
\endsong

\beginsong{Le sud}[by={Nino Ferrer}]
\textnote{Capo 4}
\beginverse
\[C]C'est un endroit qui res\[Em/B]semble à la Loui\[Am]siane, \[Am/G] \[Am/F#]
\[F]À l'Ita\[C]lie.
\[C]Il y a du linge é\[Em/B]tendu sur la ter\[Am]rasse \[Am/G] \[Am/F#]
\[F]Et c'est jo\[Em]li.
\endverse
\beginchorus
\[F]On dirait le \[C]sud, \[C7]
\[F]Le temps dure long\[C]temps, \[C7]
\[F]Et la vie sûre\[C]ment, \[Em/B] \[Am]
Plus \[Am/G]d'un million d'an\[D7]nées \[Dsus4] \[D7] \[Em] \[Dm7] \[G]
Et toujours en é\[C]té.
\endchorus
\beginverse
\[C]Y'a plein d'enfants qui se \[Em/B]roulent sur la pe\[Am]louse, \[Am/G] \[Am/F#]
\[F]Y'a plein de \[C]chiens.
\[C]Y'a même un chat, une \[Em/B]tortue, des poissons \[Am]rouges, \[Am/G] \[Am/F#]
\[F]Il ne manque \[Em]rien.
\endverse
\beginverse
\[C]Un jour ou l'autre, il \[Em/B]faudra qu'il y ait la \[Am]guerre, \[Am/G] \[Am/F#]
\[F]On le sait \[C]bien.
\[C]On aime pas ça mais on \[Em/B]ne sait pas quoi \[Am]faire, \[Am/G] \[Am/F#]
\[F]On dit : « c'est le des\[Em]tin. »
\endverse
\beginchorus
\[F]Tant pis pour le \[C]sud, \[C7]
\[F]C'était pourtant \[C]bien, \[C7]
\[F]On aurait pu \[C]vivre, \[Em/B] \[Am]
Plus \[Am/G]d'un million d'an\[D]nées \[Dm]
\[G6]Et toujours en é\[C]té.
\endchorus
\endsong

\beginsong{Love}[by={Jean-Claude Gianadda}]
\textnote{Capo 5}
\beginverse
\[Am]Il est venu chez nous, gui\[Dm]tare en bandoulière,
Venait d'on ne sait \[G]où, il parcourait la \[C]terre.
Et dans ses longs che\[E]veux, le vent semblait chan\[Am]ter,
Tout au fond de ses \[E]yeux dansait la liber\[Am]té.
\endverse
\beginchorus
\[Am]Love, c'était son \[Dm]nom, \[G]Love,
Un vaga\[C]bond qui vivait de so\[E]leil,
D'espace et de chan\[Am]sons.
\endchorus
\beginverse
\[Am]Il écoutait le vent, les \[Dm]fleurs et les rivières,
Jouait comme un en\[G]fant, parlait de la lu\[C]mière.
Il partageait ses \[E]rires, ses rêves et ses pro\[Am]jets,
Et dans chaque sou\[E]rire dansait la liber\[Am]té.
\endverse
\beginverse
\[Am]Il est parti un jour, \[Dm]nul ne sait où il est,
Au pays de l'a\[G]mour, tu peux le rencon\[C]trer.
Et dans notre mai\[E]son, il nous aura lais\[Am]sé,
Avec quelques chan\[E]sons, un peu de liber\[Am]té.
\endverse
\endsong

\beginsong{Mistral gagnant}[by={Renaud}]
\beginverse
\[Bm]À m'asseoir sur un banc, cinq minutes avec toi
Et regarder les gens, \[Em]tant qu'il y en a.
\[A7]Te parler du bon temps qu'est mort ou qui r'viendra
En serrant dans ma main \[D]tes p'tits doigts.
\[Bm]Pis donner à bouffer à des pigeons idiots,
Leur filer des coups d'pied \[Em]pour de faux.
\[A7]Et entendre ton rire qui lézarde les murs,
Qui sait surtout guérir \[D]mes blessures.
\endverse
\beginchorus
\[G]Te raconter un peu
Comment j'é\[F#7]tais, mino,
Les bombecs fabu\[Bm]leux qu'on piquait chez l'marchand.
\[G]Car en sac et Mintho,
Caramels \[F#7]à un franc,
Et les Mistral ga\[Bm]gnants.
\[A]Et les Mistral ga\[G]gnants.
\endchorus
\beginverse
\[Bm]Ah... marcher sous la pluie, cinq minutes avec toi
Et regarder la vie, \[Em]tant qu'y en a.
\[A7]Te raconter la Terre en te bouffant des yeux
Et parler de ta mère \[D]un petit peu.
\[Bm]Et sauter dans les flaques pour la faire râler,
Bousiller nos godasses \[Em]et s'marrer.
\[A7]Et entendre ton rire comme on entend la mer,
S'arrêter, r'partir \[D]en arrière.
\endverse
\beginverse
\[Bm]Ah... m'asseoir sur un banc, cinq minutes avec toi
Et regarder le soleil \[Em]qui s'en va.
\[A7]Te parler du bon temps qu'est mort et je m'en fous,
Te dire que les méchants \[D]c'est pas nous.
\[Bm]Que si moi je suis barge, ce n'est que de tes yeux,
Car ils ont l'avantage \[Em]d'être deux.
\[A7]Et entendre ton rire s'envoler aussi haut
Que s'envolent les cris \[D]des oiseaux.
\endverse
\beginchorus
\[G]Te raconter enfin
Qu'il faut ai\[F#7]mer la vie
Et l'aimer même \[Bm]si le temps est assassin
\[G]Et emporte avec lui
Les rires \[F#7]des enfants
Et les Mistral ga\[Bm]gnants.
\[A]Et les Mistral ga\[G]gnants.
\[F#7]Et les Mistral ga\[Bm]gnants.
\endchorus
\endsong

\beginsong{San Francisco}[by={Maxime Le Forestier}]
\textnote{Capo 1}
\beginverse
\[Em]C'est une maison \[G]bleue
\[Bm/F#]Adossée à la col\[C]line
On y vient à \[G]pied, on ne frappe pas
\[D/A]Ceux qui vivent là, ont jeté la \[C/G]clé.
\[Em]On se retrouve en\[G]semble
\[Bm/F#]Après des années de \[C]route
Et on vient s'as\[G]seoir, autour du repas,
\[D/A]Tout le monde est là, à cinq heures du \[A]soir.
\endverse
\beginchorus
\[Em]Quand San Fran\[G]cisco s'embrume,
\[A]Quand San Fran\[C/G]cisco s'allume,
\[D/A]San Fran\[Bm/F#]cisco,
\[Bm/F#]Où êtes-\[Em]vous,
\[G]Lizzard et \[A]Luc, Psyl\[C]via...
Atten\[Em]dez-\[Bm7]moi.
\endchorus
\beginverse
\[Em]Nageant dans le brouil\[G]lard,
\[Bm/F#]Enlacés roulant dans l'\[C]herbe,
On écoute\[G]ra, Tom à la guitare,
\[D/A]Phil à la kena, jusqu'à la nuit \[C/G]noire.
\[Em]Un autre arrive\[G]ra
\[Bm/F#]Pour nous dire des nou\[C]velles
D'un qui revien\[G]dra, dans un an ou deux,
\[D/A]Puisqu'il est heureux, on s'endormi\[A]ra.
\endverse
\beginverse
\[Em]C'est une maison \[G]bleue
\[Bm/F#]Accrochée à ma mé\[C]moire,
On y vient à \[G]pied, on ne frappe pas,
\[D/A]Ceux qui vivent là, ont jeté la \[C/G]clé.
\[Em]Peuplée de cheveux \[G]longs,
\[Bm/F#]De grands lits et de mu\[C]sique,
Peuplée de lu\[G]mière, et peuplée de fous,
\[D/A]Elle sera dernière, à rester de\[A]bout.
\endverse
\endsong

\beginsong{T'en fais pas la vie est belle}[by={Jamboree}]
\beginchorus
\[C]T'en fais pas la vie est \[G]belle
\[Am]Comme un vol d'hiron\[Em]delles
\[F]Qui s'en va et qui re\[C]vient
\[Dm]Au printemps, un beau ma\[G]tin
\endchorus
\endsong

\beginsong{Tous les cris les SOS}[by={Daniel Balavoine}]
\beginverse
\[Am7]Comme un fou va jeter à la \[F]mer
Des bou\[G]teilles vides et puis espère
Qu'on pourra \[Em]lire à tra\[F]vers.
\[Am7]S.O.S. écrit avec de l'\[F]air
Pour te \[G]dire que je me sens seul,
Je des\[Em]sine à l'encre vide un dé\[F]sert. \[G7]
\endverse
\beginchorus
Et je \[C]cours,
\[G]Je me raccroche à la \[Am]vie,
Je me \[Fmaj9]saoule avec le \[C]bruit
Des \[G]corps qui m'en\[F]tourent.
\[G]Comme des lianes nouées de \[F]tresses,
Sans com\[G]prendre la dé\[C]tresse
Des \[G]mots que j'en\[F]voie.
\endchorus
\beginverse
\[Am7]Difficile d'appeler au se\[F]cours
Quand tant de \[G]drames nous oppressent
Et les \[Em]larmes nouées de \[F]stress.
\[Am7]Étouffent un peu plus les cris d'a\[F]mour
De ceux qui \[G]sont dans la faiblesse,
Et dans un \[Em]dernier espoir dispa\[F]raissent. \[G7]
\endverse
\beginverse*
\[Dm]Tous les cris, les S.O.S.
\[Esus4]Partent dans les \[E7]airs,
Dans l'\[Am]eau, laissent une \[F]trace
Dont les \[G]écumes font la beau\[Em]té. \[F]
\[Dm]Pris dans leur vaisseau de verre,
\[Esus4]Les messages \[E]luttent,
Mais les \[Am]vagues les ra\[F]mènent
En \[G]pierres d'étoile sur les ro\[Em]chers. \[F]
\endverse*
\beginverse
\[Am7]Et j'ai ramassé les bouts de \[F]verre,
J'ai re\[G]collé tous les morceaux,
Tout était \[Em]clair comme de l'\[F]eau.
\[Am7]Contre le passé y'a rien à \[F]faire,
Il fau\[G]drait changer les héros
Dans un \[Em]monde où le plus beau reste à \[F]faire. \[G7]
\endverse
\endsong

\beginsong{Un vol d'hirondelle}
\beginverse
\[C]Comme un vol d'hiron\[F]delles
\[C]Qui s'en va et qui re\[G7]vient
\[C]Au printemps un beau ma\[F]tin
\[C]Partez là-\[G7]bas, au bout du che\[C]min
\endverse
\beginchorus
\[F]Partez, par\[C]tez
\[G7]Prenez le vent, la \[C]route et le chemin
\[F]Partez, par\[C]tez
\[G7]Vers l'horizon loin\[C]tain
\endchorus
\beginverse
\[C]Ouvrez vos \[F]ailes
\[C]Allez voir plus \[G7]loin
\[C]Portés par la \[F]brise
\[C]Jusqu'au \[G7]matin \[C]
\endverse
\endsong

\beginsong{La Blanche Hermine}[by={Gilles Servat}]
\beginverse
\[Am]J'ai rencontré ce matin
Devant la \[G]haie de mon champ
Une \[F]troupe de marins
D'ou\[Em]vriers, de pay\[Am]sans.
\endverse
\beginverse
\[Am]Où allez-vous camarades
Avec vos \[G]fusils chargés ?
Nous \[F]tendrons des embuscades,
Viens re\[Em]joindre notre ar\[Am]mée.
\endverse
\beginchorus
\[Am]La voilà la Blanche Hermine,
Vive la \[G]mouette et l'ajonc !
La \[F]voilà la Blanche Hermine,
Vive Fou\[Em]gères et Clis\[Am]son !
\endchorus
\beginverse
\[Am]Ma mie dit que c'est folie
D'aller \[G]faire la guerre aux Francs.
Moi je \[F]dis que c'est folie
D'être en\[Em]chaîné plus long\[Am]temps.
\endverse
\beginverse
\[Am]Elle aura bien de la peine
Pour é\[G]lever les enfants.
Elle \[F]aura bien de la peine
Car je \[Em]m'en vais pour long\[Am]temps.
\endverse
\beginverse
\[Am]Et si je meurs à la guerre,
Pourra-\[G]t-elle me pardonner
D'avoir \[F]préféré ma terre
À l'a\[Em]mour qu'elle me don\[Am]nait ?
\endverse
\endsong
